\documentclass[a4paper,12pt,reqno]{article}
%%%%%%%%%%%%%%%%%%%%%%%%%%%%%%%%%%%%%%%%%%%%%%%%%%%%%%%%%%%%%%%%%%%%%%%%%%%%%%%%%%%%%%%%%%%%%%%%
% 20060908 v.1.1          first public version
% 20060914 v.1.2          added a remark on the relation between B and \beta
% 20060929 v.1.3          typos + polishing for arXiv submission 
% 20060929 arxiv          arxiv version
% 20061213 cpde           submission to cpde
% 20070629 corrected      amend error on the admissible region
% 20070710 corrected      added Euler scheme
 
\usepackage{amsmath,amsfonts}
%\usepackage{mathpazo}
%\usepackage{pdfsync}
%\usepackage{showkeys}
%\oddsidemargin 0.15in  \evensidemargin 0.15in  \textwidth 6.3in
\usepackage[width=6.3in,height=24cm]{geometry} 
\DeclareMathOperator{\expect}{{\mathbb E}}
\newcommand{\eps}{\varepsilon}
\newcommand{\ott}{[0,T]}
%\newcommand{\qed}{\hfill$\square$}
%\newcommand{\proof}{{\sc Proof.}\hspace{0.3cm}}
\makeatletter
\newcommand{\eqcolon}{\mathrel{\mathord{=}\raise.2\p@\hbox{:}}}
\newcommand{\coloneq}{\mathrel{\raise.2\p@\hbox{:}\mathord{=}}}
\makeatother
\newcommand{\der}{\delta}
\newcommand{\XX}{\mathbb{X}}
\newcommand{\RR}{\mathbb{R}}
\newcommand{\TT}{\mathbb{T}}
\newcommand{\CC}{\mathcal{C}}
\newcommand{\MM}{\mathcal{M}}
\newcommand{\DD}{\mathcal{D}}
\newcommand{\CA}{\mathcal{A}}
\newcommand{\QQ}{\mathcal{Q}}
\newcommand{\LL}{\mathcal{L}}
\newcommand{\cR}{\mathcal{R}}
\newcommand{\cB}{\mathcal{B}}
\newcommand{\ZZ}{\mathbb{Z}}
\newcommand{\NN}{\mathbb{N}}
\newcommand{\OCC}{\Omega \mathcal{C}}
\newcommand{\ie}{\emph{i.e.}}
\newcommand{\norm}[1]{\lVert #1\rVert}
\newcommand{\Norm}[1]{\left\lVert #1\right\rVert}
\newcommand{\abs}[1]{\lvert #1\rvert}
\newcommand{\Abs}[1]{\left\lvert #1\right\rvert}
\newcommand{\wt}{\widetilde}
\newcommand{\wh}{\widehat}

\newcommand{\cZ}{\mathcal{Z}}
\DeclareMathOperator{\id}{\text{Id}}

\newtheorem{theorem}{Theorem}
\newtheorem{acknowledgement}[theorem]{Acknowledgement}
\newtheorem{algorithm}[theorem]{Algorithm}
\newtheorem{axiom}[theorem]{Axiom}
\newtheorem{case}[theorem]{Case}
\newtheorem{claim}[theorem]{Claim}
\newtheorem{conclusion}[theorem]{Conclusion}
\newtheorem{condition}[theorem]{Condition}
\newtheorem{conjecture}[theorem]{Conjecture}
\newtheorem{corollary}[theorem]{Corollary}
\newtheorem{criterion}[theorem]{Criterion}
\newtheorem{definition}[theorem]{Definition}
\newtheorem{example}[theorem]{Example}
\newtheorem{exercise}[theorem]{Exercise}
\newtheorem{lemma}[theorem]{Lemma}
\newtheorem{notation}[theorem]{Notation}
\newtheorem{problem}[theorem]{Problem}
\newtheorem{proposition}[theorem]{Proposition}
\newtheorem{remark}[theorem]{Remark}
\newtheorem{solution}[theorem]{Solution}
\newtheorem{summary}[theorem]{Summary}
\newenvironment{proof}[1][Proof]{\textbf{#1.} }{\ \rule{0.5em}{0.5em}}

\begin{document}
\title{Rough solutions for the\\ periodic Korteweg--de~Vries equation}
\author{
  { Massimiliano Gubinelli}              \\
%  {\small\it Dipartimento di Matematica Applicata,
%Universit\`{a} di Pisa}    \\[-0.1cm]
%  {\small\it and}          \\[-0.1cm]
{\small\it Equipe de probabilit\'es, statistique et mod\'elisation }\\[-0.1cm]
  {\small\it  Universit\'e Paris Sud}          \\[-0.1cm]
%  {\small\it Via Bonanno 25B, I-56126 Pisa, Italia.}          \\[-0.1cm]
{\small \it 91 405 Orsay cedex, France}\\[-0.1cm] 
 {\small  {\tt gubinell@math.u-psud.fr}}   
}
%\date{September 2006 (v.1.2)}
\maketitle

%\cite{*}

\begin{abstract}
The one dimensional Korteweg--de~Vries (KdV) equation on a periodic
domain and with initial condition in negative Sobolev spaces is
studied using ideas from the theory of rough paths. We discuss
convergence of Galerkin approximations, a modified Euler scheme and the presence of a random force of white-noise type in time.  
\end{abstract}

In this note we consider the Cauchy problem for the
 classical   Korteweg--de~Vries (KdV) equation:
\begin{equation}
  \label{eq:kdv-real}
\partial_t u(t,\xi) + \partial^3_\xi u(t,\xi) +
\frac12 \partial_\xi u(t,\xi)^2 = 0,\quad u(0,\xi) = u_0(\xi), \qquad
(t,\xi)\in\RR\times \TT  
\end{equation}
where  the initial condition $u_0$ belongs to some Sobolev space $H^\alpha(\TT)$
of the torus $\TT = [-\pi,\pi]$. Since the works of Bourgain~\cite{MR1215780,MR1423623} and
Kenig--Ponce--Vega~\cite{MR1329387,MR1230283,MR1086966}, it is known that eq.~(\ref{eq:kdv-real})  admits local
solutions starting from initial conditions in $H^\alpha$ for any
$\alpha >-1/2$.  Later Colliander--Keel--Staffilani--Takaoka--Tao~\cite{MR1969209} proved global well posedness for $\alpha \ge -1/2$. The existence of solutions in negative Sobolev spaces is possible due to the
regularizing effect of the dispersive linear term. This regularization
is more effective in the whole line and there global solutions exists for any $\alpha > -3/4$~\cite{MR1230283,MR1969209}. Other references on the analysis of the KdV equation are~\cite{MR1019307,MR759907}.
More recently  Kappeler--Topalov, taking advantage of the complete integrability of this model, extended the well-posedness to any $\alpha \ge -1$ using the inverse scattering method~\cite{MR2181576}.  

Inspired by the theory of \emph{rough paths} we will look for an alternative approach to the
construction of (local) solutions of eq.~(\ref{eq:kdv-real}).  

Rough paths (introduced by Lyons in~\cite{MR1654527,MR2036784}) allows the study of differential
equations driven by irregular functions. They have been applied to the
path-wise study of stochastic differential equations driven by
Brownian motion, by fractional Brownian motion of any index $H>1/4$~\cite{MR1883719}
and other stochastic processes~\cite{MR2164030,fvsubell}. Part of the theory has been reformulated in terms of the sub-Riemmannian geometry of certain Carnot groups~\cite{fvgeo}. 
In~\cite{MR2091358} we showed how to reinterpret the work of Lyons in the terms of a
cochain complex of finite increments and a related integration
theory. The key step is the introduction of a map (called
$\Lambda$) which encodes a basic fact of rough path theory.
We exploited this point of view 
 to treat stochastic partial differential
equations of evolution type~\cite{youngspde,TindelGubinelli} and to
study the initial value problem for a partial differential  equation
modeling the approximate evolution of random vortex filaments in 3d
fluids~\cite{evolution}. 

We would like to show that the concepts of the theory can be used
fruitfully for problems not related to stochastic processes.
% the analysis of some  infinite dimensional differential
%equation.
The periodic Korteweg--de Vries equation is used as a case study for our
ideas. The point of view here developed should be applicable also to KdV on
the full  real line or other dispersive semi-linear equations like the
modified KdV or the non-linear Schr\"odinger equation.  

No previous knowledge of rough paths theory is required nor any
 result on the periodic KdV is used in the following. We took care to
make the paper, as much as possible, self-contained. Since the
arguments are similar to those used in finite dimensions
in~\cite{MR2091358},  the reader can refer to this last paper to gain a wider perspective
on the technique and on the stochastic applications of rough paths. 

Our method suggests natural ways to numerically approximate
 solutions of KdV with irregular initial conditions or with random
 forcing. We describe a numerical scheme for KdV following the lines of the numerical study of SDEs driven by rough paths~\cite{davie,MR1470933,fveuler}.

\bigskip
The plan of the paper is as follows. In Sect.~\ref{sec:intro} we start by recasting the KdV
eq.~(\ref{eq:kdv-real}) in its mild form and to perform some
manipulation to motivate the finite-increment equation which we will study
in Sect.~\ref{sec:rough}  where we prove  existence and uniqueness of
local solutions. Then we prove a-priori estimates necessary
show that global rough solutions
exists for initial conditions in $L^2$ (Sect.~\ref{sec:global}).
In
Sect.~\ref{sec:galerkin} we  prove that the rough solutions are
limits of Galerkin approximations making rigorous the link between our
finite-increment problem and the classical form of the KdV
equation. Moreover in Sect.~\ref{sec:euler} we introduce a
time discretization scheme and prove its convergence.
 The
equation that we study is just one of a stack of finite-increment equations that can be
generated starting from KdV. In Sect.~\ref{sec:more} we derive the
next member of this hierarchy. Finally
Sect.~\ref{sec:stoch} addresses the problem of the presence of an
additive stochastic forcing.

We collect some longer and technical proofs in the Appendix.

\bigskip

\section{Formulation of the problem}
\label{sec:intro}
The mild form of eq.~(\ref{eq:kdv-real})  is
\begin{equation}
\label{eq:kdv-mild}
u(t,\xi) = U(t) u_0(\xi) + \frac{1}{2} \int_0^t U(t-s) \partial_\xi
u(s,\xi)^2\, ds ,\qquad t \in[0,T], \xi \in \TT  
\end{equation}
where $T>0$ is a fixed horizon and  $U(t)$ is the Airy group of isometries of any $H^\alpha$ given by the
solution of the linear part of eq.~(\ref{eq:kdv-real}). If we denote with
$\wh \varphi$ the Fourier coefficients of a function $\varphi : \TT \to
\RR$:
$
\varphi(\xi) = \sum_{k \in \ZZ} \wh \varphi(k) e^{i k \xi}
$
then the group $U(t)$ acts on $\varphi$ as
$
\wh{U(t) \varphi}(k) = e^{i k^3 t} \wh \varphi(k)
$.
Passing to Fourier variables in eq.~(\ref{eq:kdv-real}) and assuming
that the initial condition satisfy $\wh u_0(0) = 0$ we obtain the
integral equation
\begin{equation}
  \label{eq:kdv-spectral}
\wh u(t,k) = e^{i k^3 t} \wh u_0(k) + \frac{ik}2 \int_0^t e^{i k^3 (t-s)}
\sum'_{k_1} \wh u(s,k_1) \wh u(s,k-k_1) \, ds, \quad t \in [0,T], k \in \ZZ_* 
\end{equation}
where $\ZZ_* = \ZZ \backslash \{ 0 \} $ and where we use the
convention that the primed summation over $k_1 \in \ZZ$ does not include the
terms when $k_1 = 0$ or $k_1 = k$. Indeed  if
$\wh u_0(0) = 0$ then we will have $\wh u(t,0) = 0$ for any $t \ge 0$.

A last transformation of the original equation consist in using as
unknowns the variables $v(t,k) = e^{-ik^3 t} \wh u(t,k)$. The integral
equations which result is then
\begin{equation}
  \label{eq:kdv-base}
 v(t,k) =   v_0(k) + \frac{ik}2 \sum'_{k_1} \int_0^t e^{- i 3 k k_1 k_2 s}
  v(s,k_1) v(s,k_2) \, ds, \quad t \in [0,T], k \in \ZZ_* 
\end{equation}
where $k_2 = k-k_1$ and $v_0(k) = u_0(k)$. In obtaining this equation
we have used the algebraic identity $k^3 - k_1^3 - k_2^3 = 3 k k_1
k_2$ valid for any triple $k,k_1,k_2$ such that $k = k_1 + k_2$.

In the rest of the paper we will concentrate on the study of solution
to eq.~(\ref{eq:kdv-base}) with the fixed point method. Moreover since
we will always work in Fourier space we will consider $H^\alpha$ as a space
of sequences indexed by $\ZZ$ with prescribed degree of summability. 

Let us fix some notations:
we define the spaces $H^\alpha = H^\alpha(\TT)$ as the completion of the space
of  functions $\varphi:\ZZ \to\CC$ with a finite number of non-zero
values and such that $\varphi(-k) = \overline{\varphi(k)}, k \in \ZZ$ with respect to the norm
$
| \varphi |^2_\alpha \coloneq \sum_{k \in \ZZ} (1+|k|)^{2\alpha} | \varphi(k)|^2
$.
Moreover we denote $H_0^\alpha$ the subspace of $H^\alpha$ which consists of
sequences $\varphi$ such that $\varphi(0) = 0$. Note that on $H_0^\alpha$ a norm equivalent
to $|\cdot |_\alpha$ is given by
$
\varphi \mapsto (\sum_{k \in \ZZ_*} |k|^{2\alpha} |\varphi(k)|^2)^{1/2}
$.
For convenience we will usually use the latter instead that the more
usual one. The symbol $C$ in the r.h.s. of  estimates denote a positive constant which can be different from line to line. 

\subsection{Generalized integration}
In trying to solve eq.~(\ref{eq:kdv-real}) we must face two
difficulties: the presence of the $k$ factor and of the
convolution, both will come into play when trying to get
estimates for the nonlinear term in some $H^\alpha$ space. 
The only hope we have to regularize both contributions is
to exploit the oscillatory character of the phase factor $e^{-i 3 k k_1
  k_2 s}$. However it is not evident how to bound the integral
taking into account the cancellations due to this multiplier. 

To gain some insights we would like to perform a Taylor expansion of
the eq.~(\ref{eq:kdv-base}). Before  doing this we take a more
general perspective and instead of solving eq.~(\ref{eq:kdv-base}) we
will try to solve the following problem:
find a function $ v$ such
that $ v(0,k) =  v_0(k)$ and 
\begin{equation}
  \label{eq:kdv-base-2}
 v(t,k) =   v(s,k) + \frac{ik}2 \sum'_{k_1} \int_s^t e^{- i 3 k k_1 k_2 \sigma}
  v(\sigma,k_1)  v(\sigma,k_2) \, d\sigma, \quad (t,s) \in [0,T]^2, k \in \ZZ_*. 
\end{equation}
Of course if we can solve eq.~(\ref{eq:kdv-base-2}) we will have a
solution of eq.~(\ref{eq:kdv-base}). Moreover
eq.~(\ref{eq:kdv-base-2}) encodes the fact that we expect any ``piece''
of the solution to solve the same integral equation with the right
initial condition.
Now we can perform an expansion of the r.h.s of
eq.~(\ref{eq:kdv-base-2}) by substituting to each occurrence of $
v(\sigma,k)$ the corresponding expression given by the integral
constraint: a first substitution results in
\begin{equation}
  \label{eq:kdv-base-remainder}
  \begin{split}
 v(t,k) & =   v(s,k) + \frac{ik}2 \sum'_{k_1} \int_s^t e^{-i 3 k k_1 k_2 \sigma}
  d\sigma  v(s,k_1)  v(s,k_2)  + r(t,s,k)
  \end{split}
\end{equation}
where
\begin{equation*}
 % \label{eq:kdv-base-step-1}
  \begin{split}
 r(t,s,k) & = 2 \frac{ik}2 \sum'_{k_1} \int_s^t e^{-i 3 k k_1 k_2 \sigma}
  v(\sigma,k_1) \left( \frac{ik_2}2 \sum'_{k_{21}} \int_s^\sigma
   e^{- i 3 k_2 k_{21} k_{22} \sigma_1}
  v(\sigma_1,k_{21})  v(\sigma_1,k_{22}) \, d\sigma_1 \right)  \, d\sigma   
\\ &   + \frac{ik}2 \sum'_{k_1} \int_s^t e^{-i 3 k k_1 k_2 \sigma}
\left( \frac{ik_1}2 \sum'_{k_{11}} \int_s^\sigma e^{-i 3 k_1 k_{11} k_{12} \sigma_1}
  v(\sigma_1,k_{11})  v(\sigma_1,k_{12}) \, d\sigma_1 \right)
\\ &\qquad \qquad  \times
\left( \frac{ik_2}2 \sum'_{k_{21}} \int_s^\sigma e^{-i 3 k_2 k_{21} k_{22} \sigma_2}
  v(\sigma_2,k_{21})  v(\sigma_2,k_{22}) \, d\sigma_2 \right)\, d\sigma.
  \end{split}
\end{equation*}
This equation shows already an interesting pattern: the expansion
allowed us to take out from the $\sigma$-integral the unknown at the
price of introducing a correction term $r(t,s,k)$ in the form of multiple
\emph{iterated integrals}. 
 Let us pretend
for a moment that this $r$ term is   irrelevant  and let us concentrate in the analysis of the
remaining term. The oscillatory integral can now be  explicitly
evaluated:
\begin{equation}
  \label{eq:integral}
  \frac{ik}2  \int_s^t e^{-i 3 k k_1 k_2 \sigma}
  d\sigma = \frac{e^{-i 3 k k_1 k_2 s}-e^{-i 3 k k_1 k_2 t}}{6 k_1 k_2}.
\end{equation}
The $k_i$ factors at the denominator will give the regularization
which we were looking for. Introduce the symmetric bilinear operator
$X_{ts} : H^\alpha\times H^\alpha \to H^\alpha$ such that
\begin{equation}
  \label{eq:phi-op}
  X_{ts}(\varphi_1,\varphi_2)(k)  \coloneq \sum'_{k_1} \frac{e^{-i 3 k k_1 k_2
      s}-e^{-i 3 k k_1 k_2 t}}{6 k_1 k_2} \varphi_1(k_1)
  \varphi_2(k_2) .
\end{equation}
An important remark is that the family of operators $\{X_{ts},(t,s) \in
[0,T]^2\}$ satisfy the algebraic equation
\begin{equation}
  \label{eq:X-chain}
X_{ts}(\varphi_1,\varphi_2) - X_{tu}(\varphi_1,\varphi_2) -
X_{us}(\varphi_1,\varphi_2) = 0, \quad (t,u,s) \in [0,T]^3  
\end{equation}
as can be easily checked using the definition.

Define the set $\DD$ of pairs $(\gamma,\alpha)$ as $\DD = \{(\gamma,\alpha) \in [0,1/2]\times \RR  :
\gamma-\alpha < 1, 3\gamma-\alpha-3/2 < 0 \}$.
Moreover, given two Banach spaces $V,W$, denote with $\LL(V,W)$ the Banach space of bounded linear operators from $V$ to $W$ endowed with the operator norm.  

\begin{lemma}
\label{lemma:X}
For any couple
$(\gamma,\alpha) \in \DD$ and for any $\eps \ge 0$ sufficiently small,
the operator $X_{ts}$ is bounded from $H^{\alpha-\eps}\times H^{\alpha-\eps}$ to
$H^{\alpha+\eps}$ and
$
|X_{ts}|_{\LL((H^{\alpha-\eps})^2,H^{\alpha+\eps})} \le C_{\gamma,\alpha} |t-s|^\gamma
$
where the constant $C_{\gamma,\alpha}$ depends only on $\alpha$ and $\gamma$.
\end{lemma}

The proof of this lemma is postponed to Appendix~\ref{sec:reg}  .

\begin{remark}
In most of the arguments we take $\eps = 0$. A positive $\eps$ will be useful in proving convergence of Galerkin approximations below.
\end{remark}

\subsection{Computing the remainder}

Eq.~(\ref{eq:kdv-base-remainder}) has the following abstract form
\begin{equation}
  \label{eq:abs-step-1}
\der v_{ts}  = X_{ts}(v_s,v_s) + r_{ts}  
\end{equation}
where $\der v_{ts} = v_t - v_s$ and $v_t(k) = v(t,k)$ (the index
notation being more comfortable for the following discussion). The
first object on the r.h.s. is well understood thanks to
Lemma~\ref{lemma:X}, while as we have seen the $r$ term is more
complex. What is noteworthy however is that, due to the particular
structure of eq.~(\ref{eq:abs-step-1}) the $r$ term must satisfy a simple
algebraic equation. To obtain this equation, note that, for any triple
$(t,u,s) \in [0,T]^3$ we have
$
\der v_{ts} - \der v_{tu} - \der v_{us} = 0
$; substituting eq.~(\ref{eq:abs-step-1}) in this relation we get
$$
 r_{ts} -
r_{tu} - r_{us}   = 
- X_{ts}(v_s,v_s) + X_{tu}(v_u,v_u) + X_{us}(v_s,v_s).   
$$
Using eq.~(\ref{eq:X-chain}) the expression in the r.h.s. can be
simplified in
\begin{equation}
\label{eq:pre-lambda}
  \begin{split}
 r_{ts} -
r_{tu} - r_{us}  & = 
- X_{tu}(v_s,v_s) + X_{tu}(v_u,v_u)   
    \\  & = 
 X_{tu}(\der v_{us},v_s) + X_{tu}(v_u,\der v_{us})   
  \end{split}
\end{equation}
The main observation contained in the work~\cite{MR2091358} is that  sometimes this
last linear equation determines $r$ uniquely.


\subsection{The $\Lambda$-map}
To explain the conditions under which we can solve
eq.~(\ref{eq:pre-lambda}) we need some more notation. 
Given a normed vector space $(V,|\cdot|)$ introduce the vector space
$\CC_n V \subset C([0,T]^n;V)$ such that $a \in \CC_n V$ iff $ a_{t_1, \dots,t_n} = 0$ when $t_i
  = t_j$ for some $1 \le i < j \le n$. We have already introduced the
  operator $\der : \CC_1 V \to \CC_2 V$ defined as $\der f_{ts} = f_t
  - f_s$, moreover is useful to introduce
another operator $\der : \CC_2 V \to \CC_3 V$ defined on
continuous functions of \emph{two} parameters 
on a vector space $V$ as
$ \der a_{tus} = a_{ts} - a_{tu} - a_{us} $. The two operators satisfy
the relation $\der \der f = 0$ for any $f \in \CC_1 V$. Moreover if $a
\in \CC_2 V$ is such that $\der a = 0$ then  there exists $f \in \CC_1 V$ such that $ a = \der f$. Denote
$\cZ \CC_3 V = \CC_3 V \cap \text{Im} \der$.

We measure the size of elements in $\CC_n V$ for $n=2,3$  by H\"older-like norms
defined in the following way: for $f \in \CC_2 V $ let
$$
\norm{f}_{\mu} =
\sup_{s,t\in\ott}\frac{|f_{ts}|}{|t-s|^\mu},
\quad\mbox{and}\quad
\CC_1^\mu V =\{ f \in \CC_2 V;\, \norm{f}_{\mu}<\infty  \}.
$$
In the same way, for $h \in \CC_3 V$, set
\begin{eqnarray}
  \label{eq:normOCC2}
  \norm{h}_{\gamma,\rho} &=& \sup_{s,u,t\in\ott} 
\frac{|h_{tus}|}{|u-s|^\gamma |t-u|^\rho}\\
%\quad\mbox{and}\quad
\|h\|_\mu &= &
\inf\left \{\sum_i \|h_i\|_{\rho_i,\mu-\rho_i} ;\, h =
 \sum_i h_i,\, 0 < \rho_i < \mu \right\} ,\nonumber
\end{eqnarray}
where the last infimum is taken over all sequences $\{h_i \in \CC_3 V  \}$ such that $h
= \sum_i h_i$ and for all choices of the numbers $\rho_i \in (0,z)$.
Then  $\|\cdot\|_\mu$ is easily seen to be a norm on $\CC_3 V$, and we set
$
\CC_3^\mu V =\{ h\in\CC_3 V ;\, \|h\|_\mu<\infty \}
$.
Eventually,
let $\CC_n^{1+} V  = \cup_{\mu > 1} \CC_n^\mu V $ for $n =2,3$
and remark that the same kind of norms can be considered on the
spaces $\cZ \CC_3 V $, leading to the definition of the spaces
$\cZ \CC_3^\mu V $ and $\cZ \CC_3^{1+} V $. 

\vspace{0.3cm}

The following proposition is the basic result which allows the
solution of equations in the form~(\ref{eq:pre-lambda}).

\begin{proposition}[The $\Lambda$-map]
\label{prop:Lambda}
There exists a unique linear map $\Lambda: \cZ \CC^{1+}_3 V 
\to \CC_2^{1+} V $ such that 
$$
\delta \Lambda  = \id_{\cZ \CC_3 V }.
$$
Furthermore, for any $\mu > 1$, 
this map is continuous from $\cZ \CC^{\mu}_3 V $
to $\CC_2^{\mu} V $ and we have
\begin{equation}\label{ineqla}
\|\Lambda h\|_{\mu} \le \frac{1}{2^\mu-2} \|h\|_{\mu} ,\qquad h \in
\cZ \CC^{1+}_3 V . 
\end{equation}
\end{proposition}

This proposition has been first proved in~\cite{MR2091358}. A simplified proof is contained in~\cite{TindelGubinelli}.


\medskip
Using the notations just introduced, eq.~(\ref{eq:pre-lambda}) take
the form
\begin{equation}
  \label{eq:r-1}
\der r_{tus} = X_{tu}(\der v_{us},v_s) + X_{tu}(v_u, \der v_{us}) = [X(\der v,v) + X(v, \der v)]_{tus},
\qquad (t,u,s) \in [0,T]^3  
\end{equation}
where, by construction, the r.h.s. belongs to $\cZ \CC_3 H^\alpha$. If
we can prove that  it actually belongs to $ \cZ \CC^z_3 H^\alpha$
for some $z >1$, Prop.~\ref{prop:Lambda} will give us the possibility to state that the
unique solution of eq.~(\ref{eq:r-1}) in $\CC^{1+}_2 H^\alpha$ is given by
$r = \Lambda [X(\der v,v) + X(v, \der v)]$. 

\subsection{$\Lambda$ equations}

Let us come back to our initial problem. Since we aim to work in distributional spaces the rigorous meaning of eq.~(\ref{eq:kdv-base-2}) is not clear (even in a weak sense). By formal manipulations we have been able to recast the initial problem in a finite-difference equation involving the $\Lambda$ map which reads:
\begin{equation}
  \label{eq:abs-lambda}
\der v = X(v,v) + \Lambda[X(\der v, v) + X(v, \der v)]   
\end{equation}
where we used an abbreviated notation since all the terms have been
already described in detail. Note that in this equation the $\Lambda$ map has replaced the integral
 so we
will call this kind of equations: \emph{$\Lambda$-equations}. Instead
of solving the integral equation 
%(for which we do not have useful estimates) 
we would like to  solve, by fixed-point
methods, the $\Lambda$-equation~(\ref{eq:abs-lambda}). Afterwards we will show rigorously that solution to such $\Lambda$-equations gives generalized solutions to KdV.

\bigskip
  Unfortunately,
for the particular form of our $X$ operator, we are not able to show
that this equation has solutions. Recall  Lemma~\ref{lemma:X} where we showed that $X$ belongs to 
$\CC_2^\gamma \LL(H^\alpha,H^\alpha)$ for $(\gamma,\alpha) \in \DD$ so we
must expect that $\der v \in \CC_2^\gamma  H^\alpha$ since the
$\Lambda$ term will belong (at worst) to $\CC_2^{1+} H^\alpha$. But then we have
$$
X(\der v,v) + X(v,\der v) \in \CC_3^{2\gamma} H^\alpha
$$ 
so that we must require $2\gamma > 1$ in order for this term to be in
the domain of $\Lambda$. The  set $\DD$ however contains only
values of $\gamma \le 1/2$ so we are not able to study
eq.~(\ref{eq:abs-lambda}) in $\CC_1^\gamma H^\alpha$ for
$(\gamma,\alpha) \in \DD$.

This difficulty can be overcome by pursuiting further the expansion of
eq.~(\ref{eq:kdv-base-2}). By looking at
eq.~(\ref{eq:kdv-base-remainder}) we see that the next step in the
expansion will generate the following expression
\begin{equation}
  \label{eq:exp-step-2}
  \begin{split}
 v(t,k) & =   v(s,k) + \frac{ik}2 \sum'_{k_1} \int_s^t e^{-i 3 k k_1 k_2 \sigma}
  d\sigma\,  v(s,k_1)  v(s,k_2)    
\\ &   - 2 \sum'_{k_1,k_{21}}  \frac{k k_2}4  \int_s^t e^{-i 3 k k_1 k_2 \sigma}
    \int_s^\sigma
   e^{- i 3 k_2 k_{21} k_{22} \sigma_1}
 \, d\sigma_1   \, d\sigma\,  v(s,k_1)   v(s,k_{21})  v(s,k_{22})  
\\ &   + r^{(2)}(t,s,k)
  \end{split}
\end{equation}
where the new remainder $r^{(2)}$ is even more cumbersome that the
previous one and here we do not bother to write it down explicitly. In this equation a new (trilinear) operator
appear. Define $X^2 \in \CC_2 \LL((H^\alpha)^3, H^\alpha)$ as
\begin{equation}
  \label{eq:X2}
X^2_{ts}(\varphi_1,\varphi_2,\varphi_3)(k) =  
- 2 \sum'_{k_1,k_{21}}  \frac{k k_2}4  \int_s^t  \int_s^\sigma e^{-i 3
  k k_1 k_2 \sigma
- i 3 k_2 k_{21} k_{22} \sigma_1}
 \, d\sigma   d\sigma_1\,  \varphi_1(k_1)   \varphi_2(k_{21})  \varphi_3(k_{22})  
\end{equation}
and note that it satisfy the algebraic relation
\begin{equation}
  \label{eq:area-eq}
  \der X^2(\varphi_1,\varphi_2,\varphi_3)_{tus} = 2 X_{tu}(\varphi_1,X_{us}(\varphi_2,\varphi_3))
\qquad t,u,s \in [0,T]
\end{equation}
which show that $X^2$ is related to the previously defined operator
$X$ by a ``quadratic'' relationship which motivates the abuse of the superscript $2$
in the
definition of $X^2$.  



\begin{remark}
\label{rem:repr}
Define the unbounded bilinear operator $\dot X_s : (H^\alpha)^2 \to
H^\alpha$, $s\in[0,T]$ as
$$
\dot X_s(\varphi_1,\varphi_2) (k) = \frac{ik}{2}\sum'_{k_1} e^{-3k k_1 k_2
  s}\varphi_1(k_1) \varphi_2(k_2).
$$
Then for sufficiently regular $\varphi_1,\varphi_2$ we have the representations
$$
X_{ts}(\varphi_1,\varphi_2) = \int_s^t \dot
X_\sigma(\varphi_1,\varphi_2) d\sigma,
$$
and
$$
X^2_{ts}(\varphi_1,\varphi_{21},\varphi_{22}) = 2 \int_s^t \int_s^\sigma
\dot
X_\sigma(\varphi_1,\dot X_{\sigma_1}(\varphi_{21},\varphi_{22}))
d\sigma d\sigma_1,
$$
for any $(t,s) \in [0,T]^2$.
\end{remark}

The proof of the following lemma is postponed to Appendix~\ref{sec:reg}.
\begin{lemma}
\label{lemma:X2}
There exists unbounded operators $\hat X^2: (H^\alpha)^3 \to H^\alpha, \Phi:  (H^\alpha)^2 \to \mathbb{C}$ such that
$$
X^2_{ts} (\varphi_1,\varphi_2,\varphi_3) = \hat X^2_{ts} (\varphi_1,\varphi_2,\varphi_3) + \Phi_{ts}(\varphi_1,\varphi_2) \varphi_3  + \Phi_{ts}(\varphi_1,\varphi_3) \varphi_2.
$$
Furthermore   when $\alpha \ge 1/2$ we have $|\Phi_{ts}|_{\LL((H^{\alpha})^2,\mathbb{C})} \le C |t-s|$ and for any couple
$(\gamma,\alpha) \in  \DD$
we have
$
|\hat X^2_{ts}|_{\LL((H^{\alpha})^3,H^{\alpha})} \le C |t-s|^{2\gamma}
$
where the constant $C$ depends only on $\alpha$ and $\gamma$.
Then 
$
| X^2_{ts}|_{\LL((H^{\alpha})^3,H^{\alpha})} \le C'_{\gamma,\alpha} |t-s|^{2\gamma}
$
for any couple $(\gamma,\alpha) \in  \DD' = \DD \cap \{\alpha\ge 1/2\} $. Call $\DD'$ the set of \emph{allowed pairs}.
\end{lemma}


\begin{remark}
Note that the full operator $X^2$ can be controlled only in the smaller
region $\DD'$ of the $(\gamma,\alpha)$ plane. This limitation has some
relation with the low regularity ill-posedness discussed in~\cite{MR2018661}.     
\end{remark}

\bigskip
Abstractly we write eq.~(\ref{eq:exp-step-2}) as
\begin{equation}
  \label{eq:abs-step-2-ur}
\der v = X(v,v) + X^2(v,v,v) + r^{(2)}.  
\end{equation}
To infer the corresponding $\Lambda$-equation  we simply apply the $\der$ operator to
eq.~(\ref{eq:abs-step-2-ur}). After some tedious but explicit
computations we get
\begin{equation}
  \label{eq:abs-step-2-ur-rem}
  \begin{split}
\der r^{(2)}_{tus} = & 2 X_{tu}(\der v_{us}- X_{us}(v_s,v_s),v_s) 
+ X_{tu}(\der v_{us},\der v_{us}) 
\\ & + X^2_{tu}(\der v_{us},v_s,v_s) +
X^2_{tu}(v_u,\der v_{us},v_s) + X^2_{tu}( v_u,v_u,\der v_{us}).     
  \end{split}
\end{equation}

\section{Rough solutions of KdV}
\label{sec:rough}
From eq.~(\ref{eq:abs-step-2-ur-rem}) we can write down
 a second $\Lambda$-equation (beside eq.~(\ref{eq:abs-lambda})) associated to the KdV equation:
\begin{equation}
  \label{eq:abs-step-2}
\der v = X(v,v) + X^2(v,v,v) + v^\flat.  
\end{equation}
with
$$
v^\flat = \Lambda[ 2 X(\der v- X(v,v),v) 
+ X(\der v,\der v) 
 + X^2(\der v,v,v) +
X^2(v,\der v,v) + X^2( v,v,\der v)   ].
$$

An alternative form of eq.~(\ref{eq:abs-step-2}) is
\begin{equation}
  \label{eq:abs-step-2-compact}
\der v = (\id - \Lambda \der)[X(v,v)+X^2(v,v,v)]
\end{equation}
(also eq.~(\ref{eq:abs-lambda}) can be written in the same
form). While more compact this latter form is not very useful for
analysis. However it makes evident why the r.h.s. belongs to $\text{Ker} \der = \text{Im} \der$.

To give a well-defined meaning to eq.~(\ref{eq:abs-step-2}) it will be enough that the argument of $\Lambda$ belongs actually to its domain. Given an allowed pair $(\gamma,\alpha) \in \DD'$ which fixes the regularity of $X$ and $X^2$ sufficient (and natural) requirements for $v$ are
\begin{equation}
  \label{eq:requirements}
\sup_{t\in[0,T]}|v_t|_\alpha <\infty,\qquad  \der v \in \CC^\gamma_1 H^\alpha,\qquad \der v - X(v,v) \in \CC^{2\gamma}_2 H^\alpha.   
\end{equation}
Under these conditions $X(\der v- X(v,v),v) \in \CC^{3\gamma}_2 H^\alpha$ and also $X^2(\der v,v,v) +
X^2(v,\der v,v) + X^2( v,v,\der v) \in \CC^{3\gamma}_2 H^\alpha$. They are in the domain of $\Lambda$ if $3\gamma > 1$ i.e. $\gamma > 1/3$. This fixes the limiting time regularity for the $\Lambda$ eq.~(\ref{eq:abs-step-2}).

Then to define solutions of eq.~(\ref{eq:abs-step-2}) let us introduce a suitable  space to enforce the all the conditions in eq.~(\ref{eq:requirements}).
For any $0 < \eta \le \gamma$ and consider the complete metric space $\QQ_\eta$ whose elements are triples $(y,y',y^\sharp)$
where $y \in \CC_1^\eta H^\alpha$, $y' \in \CC_1^\eta H^\alpha$, $y'_0 = y_0$ and $y^\sharp \in \CC_2^{2\eta}
H^\alpha$. Additional requirement is that they satisfy the equation $\der y = X(y',y') +
y^\sharp$.  The distance $d_{\QQ,\eta}$ on $\QQ_\eta$ is defined by
$$
d_{\QQ,\eta}(y,z) = |y_0-z_0| + \|y-z\|_{\eta} + \|y'-z'\|_{\eta} + \|y^\sharp-z^{\sharp}\|_{\eta}. 
$$
for any two elements $y,z \in \QQ_\eta$. 

With abuse of notations we will denote a triple
$(y,y',y^\sharp) \in \QQ_\eta$ using simply its first component.
Moreover  we denote with $y_0$ also the constant path in $\QQ_\eta$ with value $(y_0,0,0)$. 

The main result of the paper is the following theorem.

\begin{theorem}
\label{thm:main}
Fix $\gamma \in (1/3,1/2)$ and $\alpha > 3\gamma -3/2$. Then, for any
$v_0 \in H^\alpha$ 
the $\Lambda$-equation~(\ref{eq:abs-step-2}) admits a unique solution
$v \in \QQ_\gamma$ on a small
interval $[0,T^*]$ where $T^*$ depends only on the norm of $v_0$.

If we write $v = \Theta(X,X^2,v_0)$ for the map 
$$
\Theta: \CC^\gamma_2 \LL((H^\alpha)^2,H^\alpha) \times \CC^{2\gamma}_2 \LL((H^\alpha)^3,H^\alpha) \times H^\alpha \to \QQ_\gamma
$$ which send the data to this unique local solution, then $\Theta$ is a locally Lipschitz map w.r.t. $(X,X^2,v_0) \in \CC^\gamma_2 \LL((H^\alpha)^2,H^\alpha) \times \CC^{2\gamma}_2 \LL((H^\alpha)^3,H^\alpha) \times H^\alpha$. 
\end{theorem}


The restriction on $\gamma$ and $\alpha$ ensure that the couple $(\gamma,\alpha) \in \DD'$. Then $X$ and $X^2$ are regular enough so that the proof of the theorem can
follow essentially the pattern of a similar result
in~\cite{MR2091358}.
The following is an  almost self-contained proof.

\begin{proof} Fix $\eta \in (0,\gamma)$ and
 define
a map $\Gamma: \QQ_\eta \to \QQ_\eta$ by $z = \Gamma(y)$ where
\begin{equation}
  \label{eq:gamma}
z_0 = y_0, \qquad \der z = X(y,y) + X^2(y,y,y) + z^\flat
\end{equation}
with
$$
z^\flat = \Lambda[ 2 X(y^\sharp,y) 
+ X(\der y,\der y) 
 + X^2(\der y,y,y) +
X^2(y,\der y,y) + X^2( y,y,\der y)   ].
$$
and we set $z' = y$ and $z^\sharp = X^2(y,y,y) + z^\flat$. For the map
$\Gamma$ to be well defined it is enough that $(\gamma,\alpha) \in
\DD$ and that $\eta > 1/3$. Indeed if $y \in \QQ_\eta$ then all the
terms in the argument of the $\Lambda$-map belongs to
$\CC_2^{3\eta}H^\alpha$.
 
We will proceed in three steps. First we will prove that there exists
$T_*\in(0,1]$  depending only on $|v_0|$ such that for any $T\le T_* $, $\Gamma$ maps a closed ball $B_T$ of $\QQ_{\eta}$ in
itself. Then we will prove that if $T$ is sufficiently small, $\Gamma$
is actually a contraction on this ball proving the existence  of a
fixed-point. Finally uniqueness will follow from a standard argument.


\paragraph{Boundedness.}
The key observation to prove (below) that $\Gamma$
is a contraction for $T$ sufficiently small is that $\Gamma(y)$
belongs actually to $\QQ_\gamma$, indeed $z' = y \in \CC_1^\gamma
H^\alpha$ and $z^\sharp \in \CC_2^{2\gamma} H^\alpha$. 

Let us assume that $T \le 1$.
First we will prove that
\begin{equation}
  \label{eq:bound-on-Gamma}
  d_{\QQ,\gamma}(\Gamma(y),0) = d_{\QQ,\gamma}(z,0) \le C (1+d_{\QQ,\eta}(y,0))^3.
\end{equation}
Since $3\eta > 1$ using Prop.~\ref{prop:Lambda} we have
\begin{equation}
\label{eq:flat-bound}
  \begin{split}
\|z^{\flat}\|_{3\eta} & \le C \|2 X(y^\sharp,y) 
+ X(\der y,\der y) 
 + X^2(\der y,y,y) +
X^2(y,\der y,y) + X^2( y,y,\der y) \|_{3\eta}
\\ & \le C \left[ 2\| X(y^\sharp,y) \|_{3\eta}
+ \|X(\der y,\der y)\|_{3\eta} 
\right. \\ & \qquad \left. + \|X^2(\der y,y,y)\|_{3\eta} +
\|X^2(y,\der y,y)\|_{3\eta} + \|X^2( y,y,\der y) \|_{3\eta}\right]
  \end{split}
\end{equation}
The first term is bounded by
$$
\| X(y^\sharp,y) \|_{3\eta} \le \|X\|_{\eta} \|y^\sharp\|_{2\eta} \|y\|_{\infty}
$$
where the H\"older norm $\|X\|_{\eta}$ is considered with respect to
the operator norm on $\LL((H^\alpha)^2,H^\alpha)$. Now
$\|X\|_\eta \le \|X\|_\gamma$ since we assume $\eta < \gamma$, while
$\|y^\sharp\|_{2\eta} \le d_{\QQ,\eta}(y,0)$ and finally 
remark that $\|y\|_{\infty} = |y_0|+ \|y\|_{\eta} \le
d_{\QQ,\eta}(y,0)$ since we assumed $T \le 1$. So  
$
\| X(y^\sharp,y) \|_{3\eta} \le C d_{\QQ,\eta}(y,0)^2
$.

 Each of the
other terms in eq.~(\ref{eq:flat-bound}) can be bounded similarly.
 Then 
\begin{equation}
\label{eq:flat-bound-final}
  \begin{split}
\|z^{\flat}\|_{3\eta} & \le C (1+d_{\QQ,\eta}(y,0))^3
  \end{split}
\end{equation}
This implies 
\begin{equation*}
  \begin{split}
\|z^\sharp\|_{2\gamma} & = \|X^2(y,y,y) + z^\flat\|_{2\gamma}   
%\\ &
 \le \|X^2(y,y,y)\|_{2\gamma} + \|z^\flat\|_{2\gamma}
\\ & \le C \|y\|_{\infty}^3 + \|z^\flat\|_{2\gamma}     
  \end{split}
\end{equation*}
By hypothesis we have $2\gamma < 1 < 3\eta$ so $\|z^\flat\|_{2\gamma}
\le  \|z^\flat\|_{3\eta}$ and again we obtain
$
\|z^\sharp\|_{2\gamma} \le C (1+d_{\QQ,\eta}(y,0))^3
$.
Finally we have to bound $z'$ and $z$:
\begin{equation*}
  \begin{split}
\|z'\|_{\gamma} & = \|y\|_{\gamma} \le \|X(y',y')+y^\sharp\|_{\gamma} \le
C \|y'\|_{\infty}^2 + \|y^\sharp\|_{2\gamma}
\\  & \le
C (|y_0|^2 + \|y'\|_{\gamma}^2) + \|y^\sharp\|_{2\gamma}
%\\& 
\le C (1+d_{\QQ,\eta}(y,0))^2    
  \end{split}
\end{equation*}
Similarly we can bound $z$ as $\|z\|_\gamma \le C
(1+d_{\QQ,\eta}(y,0))^3$.
Putting these bounds together we prove eq.~(\ref{eq:bound-on-Gamma}).
Let $\eps = \gamma-\eta > 0$ and observe that
\begin{equation*}
  \begin{split}
    d_{\QQ,\eta}(\Gamma(y),0) &= |z_0| +
    \|z\|_\eta+\|z'\|_{\eta}+\|z^\sharp\|_{2\eta} 
\\ & \le |z_0| +
     T^\eps\left[ \|z\|_\gamma+\|z'\|_{\gamma}+\|z^\sharp\|_{2\gamma}\right]
\\ & \le |y_0| +
     T^\eps d_{\QQ,\gamma}(\Gamma(y),0)
  \end{split}
\end{equation*}
So 
\begin{equation}
\label{eq:boundedness}
    d_{\QQ,\eta}(\Gamma(y),0) \le |y_0|+ C T^\eps (1+d_{\QQ,\eta}(y,0))^3
\end{equation}
Let $B_T = \{ y \in \QQ_{\eta} : d_{\QQ,\eta}(y,0) \le b_T, y_0 = v_0
\}$ be a closed ball of elements in $\QQ_{\eta}$ with initial
condition $v_0$. Let $b_T > 0$ be the solution of the algebraic equation
$
b_T = |v_0|+ C T^\eps (1+b_T)^3
$
which exists whenever $T \le T_*(|v_0|)$ where $T_*(|v_0|)$ is a
strictly positive time which depends only on $|v_0|$. 
If $T \le T_*(|v_0|)$ and $y \in B_T$ then by eq.~(\ref{eq:boundedness}) we have
$
d_{\QQ,\eta}(\Gamma(y),0) \le |v_0| + C T^\eps (1+b_T)^3 = b_T
$    
which implies that $\Gamma$ maps $B_T$ onto itself.

\paragraph{Contraction.}
By argument similar to those used above we can prove that if $y^{(1)},
y^{(2)}$ are two elements of $\QQ_{\eta}$ then
\begin{equation}
  \label{eq:Gamma-contraction}
  d_{\QQ,\gamma}(\Gamma(y^{(1)}),\Gamma(y^{(2)})) \le C d_{\QQ,\eta}(y^{(1)},y^{(2)}) (1+d_{\QQ,\eta}(y^{(1)},0)+d_{\QQ,\eta}(y^{(2)},0))^2 
\end{equation}
Then 
\begin{equation}
  \label{eq:Gamma-contraction-2}
   d_{\QQ,\eta}(\Gamma(y^{(1)}),\Gamma(y^{(2)})) \le
  |y^{(1)}_0-y^{(2)}_0| 
  + C T^\eps  d_{\QQ,\eta}(y^{(1)},y^{(2)}) (1+ d_{\QQ,\eta}(y^{(1)},0)+ d_{\QQ,\eta}(y^{(2)},0))^2 
\end{equation}
Taking $y_1,y_2 \in B_T$ we have
\begin{equation}
  \label{eq:Gamma-contraction-3}
  d_{\QQ,\eta}(\Gamma(y^{(1)}),\Gamma(y^{(2)})) \le
   C (1+2 b_T)^2 T^\eps d_{\QQ,\eta}(y^{(1)},y^{(2)})
\end{equation}
Let $T_{**} < T_*$ and such that  $C (1+2 b_T)^2 T_{**}^\eps < 1$, eq.~(\ref{eq:Gamma-contraction-3}) implies that
$\Gamma$ is a strict contraction on $B_{T_{**}}$ with a unique fixed-point
(in $B_{T_{**}}$).

A standard argument allows to extend the solution 
to the larger time interval $T_*$.

\paragraph{Uniqueness.} Now assume to have two solutions 
$v^{(1)},v^{(2)} \in \QQ_\eta$ of eq.~(\ref{eq:abs-step-2}), i.e. such
that $v^{(i)} = \Gamma(v^{(i)})$, $v^{(i)}_0 = v_0$ for
$i=1,2$. Eq.~(\ref{eq:Gamma-contraction}) applied to the distance between $\Gamma(v^{(1)})
$ and $\Gamma(v^{(2)})$ implies
\begin{equation*}
%  \label{eq:1}
 d_{\QQ,\gamma}(v^{(1)},v^{(2)}) \le C  d_{\QQ,\eta}(v^{(1)},v^{(2)}) (1+ d_{\QQ,\eta}(v^{(1)},0)+ d_{\QQ,\eta}(v^{(2)},0))^2   
\end{equation*}
If we denote $\QQ_{\eta,T}$ the space $\QQ_{\eta}$ considered for
functions on $[0,T]$ and with $ d_{\QQ,\eta,T}(\cdot,\cdot)$ the
corresponding distance, then we have for $S \le T$
$
 d_{\QQ,\eta,S}(y,0) \le   d_{\QQ,\eta,T}(y,0)
$ for any $y \in \QQ_{\eta,T}$. Then
\begin{equation*}
%  \label{eq:1}
 d_{\QQ,\eta,S}(v^{(1)},v^{(2)}) \le C  d_{\QQ,\eta,S}(v^{(1)},v^{(2)}) (1+ d_{\QQ,\eta,S}(v^{(1)},0)+ d_{\QQ,\eta,S}(v^{(2)},0))^2   
\end{equation*}
and so
\begin{equation*}
%  \label{eq:1}
 d_{\QQ,\eta,S}(v^{(1)},v^{(2)}) \le C S^\eps  d_{\QQ,\eta,S}(v^{(1)},v^{(2)}) (1+ d_{\QQ,\eta,S}(v^{(1)},0)+ d_{\QQ,\eta,S}(v^{(2)},0))^2   
\end{equation*}
so when $S$ is small enough that $C S^\eps
(1+ d_{\QQ,\eta,S}(v^{(1)},0)+ d_{\QQ,\eta,S}(v^{(2)},0))^2  \le 1/2  $ we
must have that $ d_{\QQ,\eta,S}(v^{(1)},v^{(2)}) = 0$ and then
$v^{(1)}_t = v^{(2)}_t$ for any $t \in [0,S]$. By repeating this
argument  we conclude that the two solution coincide in the whole interval
$[0,T]$.

\paragraph{Lipschitz continuity.}
Modulo some technicalities due to the local character of the solution, the proof of the Lipschitz continuity of the map $\Theta$ follows the line of the finite dimensional situation. See~\cite{MR2091358} for details.
\end{proof}


\subsection{$L^2$ conservation law and global solutions}
\label{sec:global}
It is well known that the KdV equation formally conserves the $L^2 =
H^0$ norm of the solution. This conservation law can be used to show
existence of global solution when the initial condition is in $H^0$. 

%\subsection{Conservative increments}
 
Denote with $\langle \cdot,\cdot \rangle$ the $L^2$ scalar product:
$
\langle \varphi,\psi \rangle = \sum_k \overline{\varphi(k)} \psi(k) =
\sum_k \varphi(-k) \psi(k)
$.

\begin{lemma}[$L^2$ differential conservation law]
For regular $\varphi_1,\varphi_2,\varphi_3$ we have
\label{lemma:basic-symmetry}
$$
\langle \varphi_1, \dot X_s(\varphi_2,\varphi_3)\rangle 
+\langle \varphi_2, \dot X_s(\varphi_1,\varphi_3)\rangle 
+\langle \varphi_3, \dot X_s(\varphi_2,\varphi_1)\rangle 
= 0, \qquad s \in[0,T]
$$  
\end{lemma}
\begin{proof}
  \begin{equation}
    \label{eq:cons-1}
\langle \varphi_1, \dot X_s(\varphi_2,\varphi_3)\rangle  =
\sum_{\substack{k,k_1,k_2\\ k=k_1+k_2}} \frac{ik}{2}e^{-i3k k_1 k_2
  s}  \varphi_{1}(-k)  \varphi_2(k_1) \varphi_3(k_2)
  \end{equation}
Make the change of variables $k \to -k$ to get
$$
- \sum_{k+k_1+k_2=0} \frac{ik}{2}e^{i3k k_1 k_2
  s}  \varphi_{1}(k)  \varphi_2(k_1) \varphi_3(k_2)
$$
Now since $k = -k_1-k_2$ this last quantity is equal to
$$
 \sum_{k+k_1+k_2=0} \frac{ik_1}{2}e^{i3k k_1 k_2
  s}  \varphi_{1}(k)  \varphi_2(k_1) \varphi_3(k_2)
+ \sum_{k+k_1+k_2=0} \frac{ik_2}{2}e^{i3k k_1 k_2
  s}  \varphi_{1}(k)  \varphi_2(k_1) \varphi_3(k_2)
$$
%In the first sum make the change of variables $k_1 \to -k_1$ and in
%the second $k_2 \to -k_2$ to get
%$$
%- \sum_{k+k_2=k_1} \frac{ik_1}{2}e^{3k k_1 k_2
%  s}  \varphi_{1}(k)  \varphi_2(-k_1) \varphi_3(k_2)
%- \sum_{k+k_1=k_2} \frac{ik_2}{2}e^{3k k_1 k_2
%  s}  \varphi_{1}(k)  \varphi_2(k_1) \varphi_3(-k_2)
%$$
%which is equal to
If 
in the first sum we make the change of variables $k_1 \to -k_1$ and in
the second $k_2 \to -k_2$ we finally get to
$
- \langle \varphi_2, \dot X_s(\varphi_1,\varphi_3)\rangle - \langle \varphi_3, \dot X_s(\varphi_2,\varphi_1)\rangle
$,
proving the statement.
\end{proof}


\begin{corollary} For $\varphi \in L^2$,
\label{cor:more-symmetry}
$$
2 \langle \varphi, X^2_{ts}(\varphi,\varphi,\varphi)\rangle
+   \langle X_{ts}(\varphi,\varphi), X_{ts}(\varphi,\varphi)\rangle = 0.
$$  
\end{corollary}
\begin{proof}
Note that Lemma~\ref{lemma:basic-symmetry} implies that
$
\langle \varphi_1 , \dot X_s(\varphi_2,\varphi_2) \rangle  =
-2 \langle \varphi_2 , \dot X_s(\varphi_1,\varphi_2) \rangle 
$
so  we can compute
  \begin{equation}
    \label{eq:cons-law-3}
    \begin{split}
  & \langle \varphi, X^2_{ts}(\varphi,\varphi,\varphi)\rangle = 
2 \int_s^t d\sigma \int_s^\sigma  d\sigma_1  
  \langle \varphi, \dot X_\sigma(\varphi,\dot
  X_{\sigma_1}(\varphi,\varphi))\rangle 
\\ & \qquad 
=- 
\int_s^t d\sigma \int_s^\sigma  d\sigma_1  
  \langle \dot
  X_{\sigma_1}(\varphi,\varphi), \dot X_\sigma(\varphi,\varphi)\rangle        
\\ & \qquad 
=- \frac12
\int_s^t d\sigma \int_s^t  d\sigma_1  
  \langle \dot
  X_{\sigma_1}(\varphi,\varphi), \dot X_\sigma(\varphi,\varphi)\rangle        
\\ & \qquad 
=- \frac12
  \langle 
  X_{ts}(\varphi,\varphi),  X_{ts}(\varphi,\varphi)\rangle        
    \end{split}
  \end{equation}
which is what we wanted to prove.
\end{proof}


\begin{theorem}[Integral conservation law]
If $v$ is a solution of eq.~(\ref{eq:abs-step-2}) in $[0,T_*]$ with initial condition $v_0 \in L^2$ then $|v_t|^2_0 = |v_0|_0^2$ for any $t \in [0,T_*]$.
\end{theorem}
\begin{proof}
We will prove that $\der \langle v,v \rangle  = 0$. Let us compute explicitly this finite increment:
\begin{equation*}
  \begin{split}
[\der \langle v,v \rangle]_{ts}  & = \langle v_t,v_t \rangle - \langle v_s,v_s \rangle
=  2 \langle \der v_{ts},v_s \rangle + \langle \der v_{ts},\der v_{ts} \rangle    
  \end{split}
\end{equation*}
Substituting in this expression the $\Lambda$-equation~(\ref{eq:abs-step-2}) we get
\begin{equation*}
  \begin{split}
[\der \langle v,v \rangle]_{ts}
& = 2  \langle X_{ts}(v_s,v_s)+X^2_{ts}(v_s,v_s,v_s) + v^\flat_{ts},v_s \rangle 
\\ & \qquad + \langle X_{ts}(v_s,v_s),X_{ts}(v_s,v_s) \rangle 
+ 2 \langle X_{ts}(v_s,v_s),v^{\sharp}_{ts} \rangle    
+  \langle v^{\sharp}_{ts},v^{\sharp}_{ts} \rangle    
  \end{split}
\end{equation*}
where we set $v^\sharp = X^2(v,v,v) +
v^\flat$. Lemma~\ref{lemma:basic-symmetry} implies that $\langle v_s,
X_{ts}(v_s,v_s) \rangle = 0$ while Corollary~\ref{cor:more-symmetry}
allows to cancel the $X^2$ term with the quadratic $X$ term. After the
cancellations the increment of the $L^2$ norm squared is  
\begin{equation*}
  \begin{split}
[\der \langle v,v \rangle]_{ts}
& = 2  \langle  v^\flat_{ts},v_s \rangle 
+ 2 \langle X_{ts}(v_s,v_s),v^{\sharp}_{ts} \rangle    
+  \langle v^{\sharp}_{ts},v^{\sharp}_{ts} \rangle    
  \end{split}
\end{equation*}
Each term  on the r.h.s. belongs at least to $\CC_2^{3\eta} \RR$ and
since $3\eta > 1$ this implies that the function $t \mapsto |v_t|^2_0$
is an H\"older function of index greater than $1$ hence it
must be constant.
\end{proof}


\begin{corollary}[$L^2$ global solutions]
If $v_0 \in L^2$ there exist a unique global solutions to the $\Lambda$-equation~(\ref{eq:abs-step-2}).  
\end{corollary}
\begin{proof}
By Thm.~\ref{thm:main} there exists a unique local solution up to a time
$T_*$ which depends only on $|v_0|_0$. Since $|v_{T_*}|_0 = |v_{0}|_0$
we can start from $T_*$ and extend uniquely this solution to the interval
$[0,2T_*]$ and then on any interval.
\end{proof}

\subsection{Galerkin approximations}
\label{sec:galerkin}
Let $P_N$ the projection on the Fourier modes $|k| \le N$. Consider the approximate KdV equation
\begin{multline}
  \label{eq:kdv-real-N}
\partial_t u^{(N)}(t,\xi) + \partial^3_\xi u^{(N)}(t,\xi) +
\frac12 P_N \partial_\xi u^{(N)}(t,\xi)^2 = 0,\\ u^{(N)}(0,\xi) = u^{(N)}_0(\xi), \qquad
(t,\xi)\in\RR\times \TT  
\end{multline}
where $u^{(N)}_0 = P_N u_0$. Is not difficult to show that the unique solution of this equation satisfy the $\Lambda$-equation
\begin{equation}
  \label{eq:abs-step-2-compact-N}
\der v^{(N)} = (\id - \Lambda \der)[X^{(N)}(v^{(N)},v^{(N)})+X^{(N),2}(v^{(N)},v^{(N)},v^{(N)})]
\end{equation}
where $v^{(N)}_t(k) = e^{-ik^3 t}\wh{u^{(N)}}(t,k)$,
$X^{(N)} = P_N X ( P_N \times P_N )$ and where the operator $X^{(N),2}$ is defined as
\begin{equation}
  \label{eq:X2-galerkin}
X^{(N),2}_{ts}(\varphi_1,\varphi_2,\varphi_3) = 2 \int_s^t d\sigma \int_s^\sigma d\sigma_1 P_N \dot X_\sigma(P_N \varphi_1,P_N \dot X_{\sigma_1}(P_N \varphi_2,P_N \varphi_3))  
\end{equation}
so that 
$$
\der X^{(N),2}(\varphi_1,\varphi_2,\varphi_3) = 2 X^{(N)}(\varphi_1,X^{(N)}(\varphi_2,\varphi_3)).
$$
These are just multi-linear operators in
a finite-dimensional space and to have convergence of the Galerkin
approximation it would be enough that both converge strongly to their
infinite-dimensional analogs $X,X^2$. Unfortunately it is not difficult to see that $X^{(N),2}$ cannot converge strongly.
 Indeed a
decomposition for   $X^{(N),2}$ analogous to that of $X^2$ described
in Lemma~\ref{lemma:X2} holds:
\begin{equation}
  \label{eq:decomp-gal}
X^{(N),2}(\varphi_1,\varphi_2,\varphi_3) = \hat X^{(N),2}(\varphi_1,\varphi_2,\varphi_3) + \Phi^{(N)}(\varphi_1,\varphi_2) P_N\varphi_3 + \Phi^{(N)}(\varphi_1,\varphi_3)P_N \varphi_2  
\end{equation}
where 
$$
\Phi^{(N)}(\varphi_1,\varphi_2) = \Phi(P_N \varphi_1,P_N \varphi_2).
$$
and we will prove in the appendix that 
\begin{lemma}
\label{lemma:gal}
Assume $(\gamma,\alpha) \in \DD$. Then as  $N \to \infty$, $X^{(N)} \to X$ strongly in $\CC_2^\gamma
\LL((H^\alpha)^2,H^\alpha)$ and $\hat X^{(N),2} \to X^2$ strongly in $\CC_2^\gamma
\LL((H^\alpha)^3,H^\alpha)$. 
\end{lemma}

Then we see that the strong convergence of $X^{(N),2}$ would imply that of the maps
$$
(\varphi_1,\varphi_2,\varphi_3) \mapsto \Phi(P_N \varphi_1,P_N \varphi_2) P_N\varphi_3 + \Phi(P_N \varphi_1,P_N \varphi_3) P_N \varphi_2
$$
and this is clearly not possible.

However the decomposition~(\ref{eq:decomp-gal}) allows us to write another $\Lambda$ equation for $v^{(N)}$
\begin{equation}
  \label{eq:abs-step-2-compact-N-bis}
\der v^{(N)} = (\id - \Lambda \der)[X^{(N)}(v^{(N)},v^{(N)})+\hat X^{(N),2}(v^{(N)},v^{(N)},v^{(N)})+ 2 \Phi(v^{(N)},v^{(N)}) v^{(N)}]
\end{equation}
note that, since $P_N v^{(N)} = v^{(N)}$, in this equation the
operator $\Phi$ appears at the place of $\Phi^{(N)}$. Of course a
similar equation holds for the solution $v$ of the full problem with
the couple $X,\hat X^2$ instead of $ X^{(N)},\hat X^{(N),2}$. 

As a consequence, the Lipschitz continuity of the solution of the
$\Lambda$-equation~(\ref{eq:abs-step-2-compact-N-bis}) w.r.t $X,\hat X^2$ and $v_0$ implies
the following convergence result.

\begin{corollary}
\label{cor:galerkin}
As $N \to \infty$ the Galerkin approximations $v^{(N)}$ converges in $\CC_1^\gamma H^\alpha$ to
the solution $v$ of the $\Lambda$-equation~(\ref{eq:abs-step-2}) up to
a strictly positive time $T^*$ which depends only on the norm of
$v_0,X,\hat X^2, \Phi$.  
\end{corollary}

This result is important in that it justifies the identification of
solutions to the $\Lambda$-equation~(\ref{eq:abs-lambda}) as
generalized solutions of the  KdV equation~(\ref{eq:kdv-real}). 



\subsection{A numerical scheme}
\label{sec:euler}
The solution described by the
$\Lambda$-equation~(\ref{eq:abs-step-2-compact}) can be approximated
by a discrete Euler-like scheme defined as follows. For any $n >0$ let $y^n_0=v_0$
and 
$$
y^n_i = X_{0, 1/n}(y^n_{i-1},y^n_{i-1}) + X^2_{0, 1/n}(y^n_{i-1},y^n_{i-1},y^n_{i-1})
$$
for $i\ge 1$. 

The combination of  this scheme with the Galerkin approximation
discussed before is an effective  numerical  approximation scheme for the
solutions of KdV with low regularity initial conditions. Indeed the next theorem  can  be combined with
Cor.~\ref{cor:galerkin} to obtain effective rates of convergence.

\begin{theorem}
Let $\Delta^n_i = y^n_i-y_{i/n}$ and let $T>0$ be the existence time of the
solution described in Thm.~\ref{thm:main}, then
\begin{equation}
  \label{eq:euler-bound}
\sup_{0\le i< j\le nT} \frac{|\Delta^n_i-\Delta^n_j|_\alpha}{|i-j|^\gamma} = O(n^{1-3\gamma}) . 
\end{equation}
\end{theorem}

\begin{proof}
Let $T$ be the existence time for the solution $v$ described in theorem~\ref{thm:main} and let $N = \lfloor nT \rfloor$.
We begin by proving some uniform bounds on the sequence
$\{y^n_i,i=0,\dots,N\}$.   Let
$$
q^n_{ij} = X_{i/n, j/n}(y^n_{i},y^n_{i}) + X^2_{i/n, j/n}(y^n_{i},y^n_{i},y^n_{i})
$$
so that $y^n_k =\sum_{i=0}^{k-1} q_{i,i+1}$. Given $0\le i<j\le N$ let
$\tau^1_0=i, \tau^1_1=j$ and define recursively  $\tau^k_l$, $k >
0$, $l =0,\dots,2^k$ such that $\tau^{k+1}_{2l} = \tau^k_l$  for
$l=0,\dots,2^k$ and 
$\tau^{k+1}_{2l+1}<\lfloor (\tau^{k}_{l}+\tau^{k}_{l+1})/2\rfloor$ for
$l=0,\dots,2^k-1$. 
Then we have $|\tau^k_{l+1}-\tau^k_l|\le 1 \vee (j-i)/2^k$
and
$
y^n_j-y^n_i =\sum_{l=0}^{2^K} q^n_{\tau^k_l,\tau^k_{l+1}} 
$
 where $K$ is such that
$(j-i)/2^K \le 1$. Using the triangular array $\tau^k_l$ we 
rewrite the above expression as a telescopic sum:
\begin{equation*}
  \begin{split}
y^n_j-y^n_i & = q^n_{ij} + \sum_{k=1}^K \sum_{l=0}^{2^{K-1}} (q^n_{\tau^k_{2l},\tau^k_{2l+2}}-q^n_{\tau^k_{2l},\tau^k_{2l+1}}-q^n_{\tau^k_{2l+1},\tau^k_{2l+2}}) 
  \\ & = q^n_{ij} + \sum_{k=1}^K \sum_{l=0}^{2^{k-1}} (\der q^n)_{\tau^k_{2l},\tau^k_{2l+1},\tau^k_{2l+2}}   
  \end{split}
\end{equation*}
Up to $T$, the solution $v$ satisfy the equation
$
\der v = X(v,v) + X^2(v,v,v) + v^\flat
$
where $\sup_{t\in[0,T]}|v_t|_\alpha+\|v^\flat\|_{3\gamma}\le C$ so
\begin{equation*}
v_{j/n}-v_{i/n} = \sum_{l=i}^{j-1} p_{l,l+1} + \sum_{l=i}^{j-1} v^\flat_{l/n,(l+1)/n}
\end{equation*}
where $p_{ij}=X_{i/n,j/n}(v_{i/n},v_{i/n})+X^2_{i/n,j/n}(v_{i/n},v_{i/n},v_{i/n})$.
Then
\begin{equation*}
\Delta^n_{j}-\Delta^n_{i} =  q^n_{ij}-p_{ij}  + \sum_{k=1}^K \sum_{l=0}^{2^{k-1}} (\der
q^n- \der p)_{\tau^k_{2l},\tau^k_{2l+1},\tau^k_{2l+2}}  - r_{ij}
\end{equation*}
where
$
r_{ij}=  \sum_{l=i}^{j-1} v^\flat_{l/n,(l+1)/n}
$.
This last term is readily estimated by
$$
|r_{ij}| \le \sum_{l=i}^{j-1} \|v^\flat\|_{3\gamma} n^{-3\gamma}
 \le C \left(\frac{j-i}{ n}\right) n^{1-3\gamma}
$$
uniformly in $n$.
Let 
$$
M^n_\ell = \sup_{0\le i < j \le \ell} \left(\frac{j-i}{ n}\right)^{-1} |\Delta^n_{j}-\Delta^n_{i} -  q^n_{ij}+p_{ij}+r_{ij}|
$$
we want to show that $M^n_N \le A n^{1-3\gamma}$ uniformly in $n$ for some constant
$A$ depending only on the data of the problem: this will imply the
statement of the theorem since then
$$
|\Delta^n_{i}| \le |q^n_{0i}-p_{0i}| + |r_{0i}| + A n^{1-3\gamma}
\left(\frac{i}{ n}\right)  \le |r_{0i}| + A n^{1-3\gamma}
\left(\frac{i}{ n}\right) \le Cn^{1-3\gamma}
$$
for any $i \le N$ and
$$
|\Delta^n_{j}-\Delta^n_{i}| \le |q^n_{ij}-p^{ij}| + |r_{ij}| + A n^{1-3\gamma}
\left(\frac{j-i}{ n}\right)  \le C n^{1-3\gamma} \left(\frac{j-i}{ n}\right)^\gamma
$$
for any $0\le i<j\le N$.


 We proceed by
induction on $\ell$. For $\ell=1$ the statement is clearly true since
$\Delta^n_{1}-\Delta^n_{0} -  q^n_{0,1}+p_{0,1}+r_{0,1}=0$, moreover
for the same reason we have, for all $l$ that
$\Delta^n_{l+1}-\Delta^n_{l} -  q^n_{l,l+1}+p_{l,l+1}+r_{l,l+1}=0$.
 Assume then
that $M^n_{\ell-1} \le A$ for some $\ell> 0$. The basic observation is that when $|i-j|\le \ell$
the sums
$$
\Delta^n_{j}-\Delta^n_{i} -  q^n_{ij}+p_{ij}+r_{ij} = \sum_{k=1}^K \sum_{l=0}^{2^{k-1}} (\der
q^n- \der p)_{\tau^k_{2l},\tau^k_{2l+1},\tau^k_{2l+2}}
$$
can be estimated in terms of $M^n_{\ell-1}$ and various norms of
$v$,$X$,$X^2$ much like in the proof of Thm.~\ref{thm:main}. The bound
has the form
$$
|(\der
q^n- \der p)_{\tau^k_{2l},\tau^k_{2l+1},\tau^k_{2l+2}}| \le C
\left(\frac{j-i}{2^k n}\right)^{3\gamma} (1+M^n_{\ell-1})^3 (M^n_{\ell-1}+n^{1-3\gamma})
$$
where $C=C(v_0,X,X^2)$ and where the factor $n^{1-3\gamma}$ is due to
the previous estimate on $r_{ij}$. Then since $3\gamma > 1$ and
$$
\sum_{k=1}^K \sum_{l=0}^{2^{k-1}}
\left(\frac{j-i}{2^k n}\right)^{3\gamma}  \le \left(\frac{j-i}{
    n}\right)^{3\gamma} \sum_{k=1}^\infty 2^{k(1-3\gamma)} \le C  \left(\frac{j-i}{
    n}\right)^{3\gamma} \le C \left(\frac{\ell}{
    n}\right)^{3\gamma-1}  \left(\frac{j-i}{
    n}\right) 
$$
we get
$$
M^n_\ell \le C
(1+M^n_{\ell-1})^4 \sum_{k=1}^K \sum_{l=0}^{2^{k-1}}
\left(\frac{j-i}{2^k n}\right)^{3\gamma} \le C (1+M^n_{\ell-1})^3 (M^n_{\ell-1}+n^{1-3\gamma}) \left(\frac{\ell}{
    n}\right)^{3\gamma-1} 
$$
Let $m^n_\ell = n^{3\gamma-1} M^{n}_\ell$, then, for $n$ large enough
$m^n_\ell \le F(m^n_{\ell-1})$ where
$F$ is the increasing map
$$
\RR_+ \ni m \mapsto F(m)= C (1+m)^4 \left(\frac{\ell}{
    n}\right)^{3\gamma-1} \in \RR_+
$$
which, for $\ell/n$ small enough
has a unique attracting fix-point under iteration starting from $0$. 
 In particular the
iterations stay bounded and if we  set $x_{i+1}=F(x_i)$, $x_0=0$ we
have $A = \sup_i F(x_i)  <
\infty$ and $m^n_\ell \le x_\ell \le A$. By repeating this argument it is easy to prove that the
bound holds for all $\ell \le nT$, i.e. in the whole existence
interval found in Thm.~\ref{thm:main}.
\end{proof}

\begin{remark}
With a bit more of work it is possible to prove the existence of
the solution stated in Thm.~\ref{thm:main} using directly the discrete
approximation as done by Davie~\cite{davie} for rough
differential equations.
\end{remark}


\subsection{Higher order $\Lambda$ equations}
\label{sec:more}
Further expansion of eq.~(\ref{eq:exp-step-2}) generate a hierarchy of
$\Lambda$-equations for KdV. The next one is given by
\begin{equation}
  \label{eq:hier-next}
  \der v = (\id-\Lambda\der) [X(v^{\times 2}) + X^2 (v^{\times 3}) + X^{3a}(v^{\times 4}) + X^{3b}(v^{\times 4})]
\end{equation}
where $v^{\times n} = (v,\dots,v)$ ($n$ times) and where $X^{3a},X^{3b}$ are operators increments respectively defined as
\begin{equation}
  \label{eq:X3a}
  X^{3a}_{ts}(\varphi_1,\varphi_2,\varphi_3,\varphi_4) = \int_s^t
  d\sigma_1 \dot
  X_{\sigma_1}(\varphi_1,\int_s^{\sigma_1}d\sigma_2 \dot X_{\sigma_2}(\varphi_2,
\int_s^{\sigma_2}d\sigma_3 \dot X_{\sigma_3}(\varphi_3,\varphi_4)))
\end{equation}
and
\begin{equation}
  \label{eq:X3b}
  X^{3b}_{ts}(\varphi_1,\varphi_2,\varphi_3,\varphi_4) = \int_s^t
  d\sigma_1 \dot
  X_{\sigma_1}(\int_s^{\sigma_1}d\sigma_2 \dot
  X_{\sigma_2}(\varphi_1,\varphi_2),
\int_s^{\sigma_1}d\sigma_3 \dot
  X_{\sigma_3}(\varphi_3,\varphi_4))
\end{equation}
which satisfy the following relations with $X$ and $X^2$:
\begin{equation}
  \label{eq:deltaX3a}
  \der X^{3a}(\varphi_1,\varphi_2,\varphi_3,\varphi_4) =
  X(\varphi_1,X^2(\varphi_2,\varphi_3,\varphi_4)) + X^2(\varphi_1,\varphi_2,X(\varphi_3,\varphi_4))
\end{equation}
and
\begin{equation}
  \label{eq:deltaX3b}
  \begin{split}
  \der X^{3b}(\varphi_1,\varphi_2,\varphi_3,\varphi_4) & =
  X(X(\varphi_1,\varphi_2),X(\varphi_3,\varphi_4)) +
  X^2(\varphi_3,\varphi_4,X(\varphi_1,\varphi_2))
\\ & \qquad  + X^2(\varphi_1,\varphi_2,X(\varphi_3,\varphi_4))    
  \end{split}
\end{equation}

As we report elsewhere~\cite{trees}  the Hopf algebra of rooted trees is
the natural language to describe this hierarchy of equations and the
algebraic relations between the various operators. In this special
case however these relations can be easily checked by direct 
 computations. 

Using Lemmas~\ref{lemma:X} and~\ref{lemma:X2} we can show that the r.h.s. of
the eqns.~(\ref{eq:deltaX3a}) and~(\ref{eq:deltaX3b}) belongs to the
domain of $\Lambda$ and so the equations can be used to express
$X^{3a},X^{3b}$ in function of $X,X^2$ and prove that

\begin{corollary}
For $(\gamma,\alpha) \in\DD'$ with $\gamma>1/3$ we have
$X^{3a},X^{3b} \in \CC_2^\gamma \LL((H^{\alpha})^4,H^{\alpha}).$ 
\end{corollary}




\section{Additive stochastic forcing}
\label{sec:stoch}
As another application of this approach we would like to discuss the presence of an additive random force in the KdV eq.~(\ref{eq:kdv-real}):
\begin{equation}
  \label{eq:kdv-real-stoch}
\partial_t u(t,\xi) + \partial^3_\xi u(t,\xi) +
\frac12 \partial_\xi u(t,\xi)^2 = \Phi \partial_t  \partial_\xi B(t,\xi)
\end{equation}
where $\Phi$ is a bounded linear operator on $H^0$ and
$\partial_t \partial_\xi B$ a white noise on $\RR\times\TT$.
We assume that $\Phi$ is diagonal in Fourier space: $\Phi e_k =
\lambda_k e_k$ where $\{e_k \}_{k\in\ZZ}$ is the orthonormal basis
$e_k(\xi) = e^{ik\xi}/\sqrt{2\pi}$ and that $\lambda_0 = 0$.  In this
way the noise does not affect the zero mode.

For the rest of this section  fix and allowed pair $(\gamma,\alpha)
\in \DD'$ such that $\gamma \in (1/3,1/2)$ and assume that
\begin{equation}
  \label{eq:lambda-bound}
\sum_{q \in \ZZ_*}  |q|^{2\alpha} \lambda_q^2 < \infty.
\end{equation}
which implies that $\Phi$ is  Hilbert-Schmidt  from $L^2$ to
$H^\alpha$.

The transformed integral equation (analogous to
eq.~(\ref{eq:kdv-base})) associated to~(\ref{eq:kdv-real-stoch}) is
\begin{equation}
  \label{eq:kdv-base-stoch}
 v(t,k) =   v_0(k) + \lambda_k \beta^k_t +  \int_0^t \dot X_s(
  v(s), v(s))(k) \, ds, \quad t \in [0,T], k \in \ZZ_* 
\end{equation}
where $\{\beta^k_\cdot\}_{k \in \ZZ_*}$ is a family of 
complex-valued, zero mean, Brownian motions such that $\beta^{-k} =
\overline{\beta^k}$ and with covariance
$\expect[\overline \beta^k_t \beta^q_s] = \delta_{k,q} (t\wedge s)$. 
The relation between the initial noise $B$ and the family $\{\beta^k_\cdot\}_{k \in \ZZ_*}$ is given by 
$
\beta^k_t = \langle e_k, \int_0^t U(-s) \partial_\xi \partial_s B(s,\cdot) ds \rangle_{H^0} 
$.

Eq.~(\ref{eq:kdv-base-stoch}) can be expanded in the same
way as we have done before and the first
interesting $\Lambda$-equation which appears is the following:
\begin{equation}
  \label{eq:lambda-eq-stoch}
  \der v = (\id-\Lambda\der)[X(v,v) + \der w + X^{2}(v,v,v) + X^{w}(v)].
\end{equation}
Here $\der w(k)_{ts} = \lambda_k (\beta^k_t-\beta^k_s)$ and  the
random operator $X^w_{ts}: H^\alpha\to H^\alpha$ is given by
\begin{equation}
  \label{eq:Xw}
  X^w_{ts}(\varphi) = \int_s^t d\sigma \dot X_\sigma(\varphi, \der
  w_{\sigma s})
\end{equation}
and satisfy the equation
\begin{equation}
  \label{eq:deltaXw}
  \der X^w_{tus} = X_{tu}(\varphi,\der w_{us}).
\end{equation}

From condition~(\ref{eq:lambda-bound}) we have that $w$ is a.s. a Brownian
motion with values in $H^{\alpha}$ and then by classical results that $\der w \in
\CC^{\rho}_2 H^\alpha$ for any $\rho < 1/2$ and a-fortiori $\der w \in
\CC^{\gamma}_2 H^\alpha$ by choosing $\rho \in [\gamma,1/2)$. 

To prove that a sufficiently regular version of the Gaussian stochastic process $X^w$ exists
we will use a generalization of the classic Garsia-Rodemich-Rumsey lemma
which has been proved in~\cite{MR2091358}.
\begin{lemma}
\label{lemma:besov}
For any $\theta > 0$ and $p \ge 1$,  there exists a constant
$C$ such that for any $R \in \CC_{2} V$ ($(V,|\cdot|)$ some Banach space), we have 
\begin{equation}
\label{eq:generalboundxx}
\|R\|_{\theta} \le C (U_{\theta+2/p,p}(R)+\|\der R\|_{\theta}),
\end{equation}
where
\begin{equation*}
 U_{\theta,p}(R) = \left[ \int_{[0,T]^2}
 \left(\frac{|R_{t s}|}{|t-s|^\theta}\right)^p dt ds \right]^{1/p}.
\end{equation*}
\end{lemma}

We would like that $X^w$ behaves not worser than $X^2$. 
This is indeed the case as stated in the next lemma.

\begin{lemma}
\label{lemma:Xw}
Under condition~(\ref{eq:lambda-bound})
we have $X^w \in \CC^{2\gamma}_2 \LL(H^\alpha,H^\alpha)$ a.s..
\end{lemma}
\begin{proof}
After an integration by parts, $X^w$ can be rewritten as
\begin{equation}
  \label{eq:Xw-2}
  X^w_{ts}(\varphi) = X_{ts}(\varphi,\der w_{ts}) - \int_s^t X_{\sigma
  s}(\varphi,dw_\sigma) = X_{ts}(\varphi,\der w_{ts}) + I_{ts} 
\end{equation}
The first term in the r.h.s. belongs to $\CC_2^{2\gamma}
\LL(H^\alpha,H^\alpha)$ path-wise:
$$
|X_{ts}(\varphi,\der w_{ts})|_{\alpha} \le
|X_{ts}|_{\LL(H^\alpha\times H^\alpha,H^\alpha)} |\der w_{ts}|_\alpha
|\varphi|_\alpha \le 
\|X\|_\gamma \|\der
w\|_\gamma |\varphi|_\alpha |t-s|^{2\gamma}.
$$
Let us estimate the Wiener integral $I_{ts}$. Take $\gamma_1 < \gamma$ so that $(\gamma_1,\alpha)$ is still an allowed pair,
\begin{equation*}
  \begin{split}
\expect \left|I_{ts}\right|^2_\alpha & = 
\sum_q \lambda_q^2 \int_s^t |X_{\sigma s}(e_q,\varphi)|_\alpha^2
d\sigma
\\ & \le C \sum_q \lambda_q^2 \int_s^t |\sigma-s|^{2\gamma_1} d\sigma \sum_k
|k|^{2\alpha}
\left|\sum_{k_1} |k|^{\gamma_1} |k_1|^{\gamma_1-1} |k_2|^{\gamma_1-\alpha-1}
  e_q(k_1) \psi(k_2)\right|^2    
\\ & \le C \sum_q \lambda_q^2 \int_s^t |\sigma-s|^{2\gamma_1} d\sigma \sum_k
|k|^{2\alpha+2\gamma_1} |q|^{2\gamma_1-2} |k-q|^{2\gamma_1-2\alpha-2}
  |\psi(k-q)|^2    
  \end{split}
\end{equation*}
where we used a bound for the kernel of $X$ which has been introduced
in the proof of Lemma~\ref{lemma:X} and where $\psi(k) =
|k|^{\alpha}\varphi(k)$. 
The function $k \mapsto
|k|^{2\alpha+2\gamma_1} |k-q|^{2\gamma_1-2\alpha-2}$ is bounded by $C
|q|^{2\alpha+2\gamma_1}$ so 
\begin{equation*}
  \begin{split}
\expect \left|I_{ts}\right|^2_\alpha &  \le C \sum_q \lambda_q^2  |q|^{4\gamma_1+2\alpha-2} 
|t-s|^{1+2\gamma_1} |\psi|_0^2
  \end{split}
\end{equation*}
Choosing $\gamma_1 = 2\gamma-1/2+\eps/4\le \gamma$ (with $\eps>0$ small enough) we obtain
\begin{equation*}
  \begin{split}
\expect \left|I_{ts}\right|^2_\alpha &  \le C \sum_q \lambda_q^2  |q|^{8\gamma+2\alpha-4+\eps} 
|t-s|^{4\gamma+\eps/2} |\varphi|_\alpha^2  \le C |t-s|^{4\gamma+\eps/2} |\varphi|_\alpha^2
  \end{split}
\end{equation*}
The $q$ sum is bounded
since we can choose $\eps>0$ small enough so that $8\gamma-4+\eps \le
0$ which allows us to bound $|q|^{8\gamma+2\alpha-4+\eps}  \le
|q|^{2\alpha}$ and then conclude using
condition~(\ref{eq:lambda-bound}) on the summability of $\lambda_q$.

Since the Wiener integral $I_{ts}$ is a Gaussian r.v.  the $L^2$ bound is
enough to obtain an $L^p$ bound for any $p\ge 2$:
\begin{equation*}
  \begin{split}
\expect \left|I_{ts}\right|^p_\alpha &  \le C_p  |\varphi|_\alpha^p |t-s|^{2p\gamma+p \eps /4}
  \end{split}
\end{equation*}
From this bound we deduce that, choosing $p$ sufficiently large, we
have $U_{2\gamma+2/p,p}(I) < \infty$ a.s.. Moreover $\der I_{tus}(\varphi)
= X_{tu}(\varphi,\der w_{us})$ so that $\der I$ can be bounded
path-wise in $\CC_3^{2\gamma} \LL(H^{\alpha},H^{\alpha})$.  
Then Lemma~\ref{lemma:besov} implies that 
$
\|I\|_{2\gamma} \le C(U_{2\gamma+2/p,p}(I)+\|\der I\|_{2\gamma}) < \infty
$ a.s.
ending the proof.
\end{proof}

Modifying a bit the proof of Thm.~\ref{thm:main} is not
difficult to prove the following.

\begin{theorem}
\label{th:eq-stoch}
Eq.~(\ref{eq:lambda-eq-stoch}) has a unique local solution in $\CC^\gamma H^\alpha$ for any initial
condition in $H^\alpha$.  
\end{theorem}

Since we can take $\gamma$ arbitrarily close to $1/3$ we obtain
solutions for noises with values in $H^\alpha$ for any  $\alpha > -1/2$.
In this way we essentially cover the result of
De~Bouard-Debussche-Tsutsumi~\cite{MR2111917}. Their approach consist
in modifying Bourgain's method to handle  Besov spaces in order to compensate
for the insufficient  Sobolev time regularity of Brownian motion.


\section*{Acknowledgment}
I would like to thank A. Debussche which  delivered a series of interesting
lectures on stochastic dispersive equations  during a 2006
semester on Stochastic Analysis at Centro de Giorgi, Pisa. They
constituted the motivation for the investigations reported in this
note. I'm also greatly indebited with J.Colliander for  some
interesting remarks which have been instrumental in  discovering an error in  an earlier version of the
paper.  

%%%%%%%%%%%%%%%%%%%%%%%%%%%
%%%%%%%%%%%%%%%%%%%%%%%%%%%
%%%%%%%%%%%%%%%%%%%%%%%%%%%

\appendix
\section{Regularity of $X$ and $X^2$}
\label{sec:reg}


Some elementary calculus results needed in the proofs of this appendix are the
subject of the next lemma.

\begin{lemma}
\label{lemma:sums}
If $\zeta > 1$ 
$$
\sum_{n\in \ZZ \backslash \{0, m \}}
    |n|^{-\zeta} |n-m|^{-\zeta} \le C_\zeta |m|^{-\zeta}
$$
and for any $\theta < 2\zeta-1$
$$
\sum_{n\in \ZZ \backslash \{0,\pm m \}}
    |n+m|^{-\zeta} |n|^{\theta} |n-m|^{-\zeta} \le C_{\zeta,\theta} |m|^{\theta-\zeta}
$$  
\end{lemma}

\subsection{Proof of Lemma~\ref{lemma:X}}

\begin{proof}[]
We must prove that, for $(\gamma,\alpha) \in \DD$ there exists $\eps > 0$ such that
$$
\left[
\sum_{k \in \ZZ_*} |k|^{2\alpha+2\eps} \left|\sum'_{k_1} \frac{e^{-i 3 k k_1 k_2
      s}-e^{-i 3 k k_1 k_2 t}}{6 k_1 k_2} \varphi_1(k_1)\varphi_2(k_2)
\right|^2 \right]^{1/2}\le C |t-s|^\gamma |\varphi_1|_{\alpha-\eps} |\varphi_2|_{\alpha-\eps}
$$
To begin we bound
$$
\left| \frac{e^{-i 3 k k_1 k_2
      s}-e^{-i 3 k k_1 k_2 t}}{6 k_1 k_2}\right| \le C |t-s|^\gamma
|k|^\gamma |k_1|^{\gamma-1} |k_2|^{\gamma-1}
$$
for any $\gamma \in[0,1]$.
We use the conventions that, if $\beta < 0$ then  $|0|^\beta =
0$ (with integer argument). The constant in this bound is independent on anything else.
By duality it will be enough to prove that, for any $\rho,\psi_1,\psi_2 \in
H^0_0$ we have
\begin{equation}
  \label{eq:dual-bound-1}
\sum_{k, k_1}  |k|^{\alpha+\eps+\gamma} |k_1 k_2|^{\gamma-\alpha+\eps-1} 
|\rho(k)| |\psi_1(k_1)| |\psi_2(k_2)| \le C |\rho|_0 |\psi_1|_{0} |\psi_2|_{0}.  
\end{equation}
where we understand that summations are done over $\ZZ$ excluding
values for which $k_1 = 0$ or $k = 0$ or $k_2 = 0$.
 
To prove eq.~(\ref{eq:dual-bound-1}) for any couple
$(\gamma,\alpha)\in \DD$ we proceed in three steps: we will
prove first that it holds for $(\gamma,\alpha)$ in two regions
$\DD_1,\DD_2 \subset \RR^2 $ and then by interpolation we will
extend its validity to the convex hull of $\DD_1 \cup
\DD_2$ which will result to be exactly $\DD$. 

\paragraph{First region.}
Let $\DD_1 = \{(\gamma,\alpha) : \gamma-\alpha < 1/2, \gamma
\in [0,1/2) \}$ and  in this paragraph fix $(\gamma,\alpha) \in \DD_1$.
Applying  Cauchy-Schwarz inequality w.r.t the $k_1$ summation on the
l.h.s. of
eq.~(\ref{eq:dual-bound-1})  we obtain
\begin{equation}
\label{eq:dual-bound-2}
  \sum_{k}  |k|^{\alpha+\gamma+\eps} |\rho(k)|  \left[\sum_{k_1}
    |k_1 k_2|^{2\gamma-2\alpha+2\eps-2}  \right]^{1/2}
  \left[\sum_{k_1 } |\psi_1(k_1)|^2 |\psi_2(k_2)|^2 \right]^{1/2} 
\end{equation}
Observe that, since $2\gamma-2\alpha-1 < 0$, 
Lemma~\ref{lemma:sums} allows us to bound the sum over $k_1$ as
$$
\sum_{k_1}
    |k_1 k_2|^{2\gamma-2\alpha+2\eps-2} \le C |k|^{2\gamma-2\alpha+2\eps-2}
$$
for sufficiently small $\eps > 0$.
Injecting this bound in eq.~(\ref{eq:dual-bound-2}) we find
\begin{equation}
\label{eq:dual-bound-2bis}
 C  \sum_{k}  |k|^{2\gamma-1+4\eps} |\rho(k)|   \left[\sum_{k_1 }
    |\psi_1(k_1)|^2 |\psi_2(k_2)|^2 \right]^{1/2} 
\end{equation}
Given that $\gamma < 1/2$ we have $|k|^{2\gamma-1+4\eps} \le 1$ if we choose $\eps$ small enough. Taking this
into account in the last bound and applying Cauchy-Schwarz w.r.t to
the $k$ summation we can bound our initial expression by
$ C  |\rho|_0   [\sum_{k, k_1 }
    |\psi_1(k_1)|^2 |\psi_2(k_2)|^2 ]^{1/2} $. But now observe that a
    change of variables is enough to estimate this last double
    summation:
$$
\sum_{k, k_1 }
    |\psi_1(k_1)|^2 |\psi_2(k_2)|^2 = \sum_{k_2  k_1 }
    |\psi_1(k_1)|^2 |\psi_2(k_2)|^2 = |\psi_1|_0^2 |\psi_2|_0^2
$$
so that we have proven eq.~(\ref{eq:dual-bound-1}) in the case of
$(\gamma,\alpha) \in \DD_1$.

\paragraph{Second region.}
Let $\DD_2 = \{(\gamma,\alpha)\in\RR^2 : \gamma \in[ 0,1/2),  \gamma-\alpha < 1, \gamma
+\alpha < -1/2 \}$ and fix $(\gamma,\alpha) \in \DD_2$.
Apply Cauchy-Schwartz to the l.h.s. of eq.~(\ref{eq:dual-bound-1}) to obtain
$$
\left[\sum_k |k|^{2\gamma+2\alpha+2\eps}\right]^{1/2}
\left[\sum_k |\rho(k)|^2 \left(\sum_{k_1} 
    |k_1 k_2|^{\gamma-\alpha-1+2\eps} |\psi_1(k_1)|  |\psi_2(k_2)| \right)^2\right]^{1/2}
$$
Choose $\eps$ small enough so that $\gamma+\alpha+\eps < -1/2$ and $\gamma-\alpha+\eps \le 1$. Then these two inequalities imply that the first sum over $k$ is bounded and that we can discard the powers of $k_1$ and $k_2$ in the second factor. Finally we get
$$
C
\left[\sum_k |\rho(k)|^2 \left(\sum_{k_1} |\psi_1(k_1)|  |\psi_2(k_2)| \right)^2\right]^{1/2}
$$
Now observe that $ \sum_{k_1} |\psi_1(k_1)|  |\psi_2(k_2)| \le
|\psi_1|_0 |\psi_2|_0$ and then we can bound the last display by
$
|\rho|_0 |\psi_1|_0 |\psi_2|_0
$, proving eq.~(\ref{eq:dual-bound-1}) in the region $\DD_2$.

\paragraph{Interpolation.}
Let $(\gamma_1,\alpha_1) \in \DD_1 \cup \DD_2$ and  $(\gamma_2,\alpha_2)
\in  \DD_1 \cup \DD_2$, consider for $t\in[0,1]$,  $\gamma(t)=t \gamma_1 +
(1-t) \gamma_2$,  $\alpha(t)=t \alpha_1 +
(1-t) \alpha_2$. Moreover let $\eps$ be the smaller of the two $\eps$ related to the couple $(\gamma_1,\alpha_1)$ and $(\gamma_2,\alpha_2)$.
Let
$$
A(\gamma,\alpha) = \sum_{k, k_1}  |k|^{\alpha+\eps+\gamma} |k_1|^{\gamma-\alpha+\eps-1} |k_2|^{\gamma-\alpha+\eps-1}
|\rho(k)| |\psi_1(k_1)| |\psi_2(k_2)|.
$$ 
By using H\"older inequality with conjugate exponents $p = 1/t$ and $q
= 1/(1-t)$ we have
$$
A(\gamma(t),\alpha(t)) \le A(\gamma_1,\alpha_1)^{1/p} A(\gamma_2,\alpha_2)^{1/q}
$$
so since we have
$
A(\gamma_i,\alpha_i) \le C_{\gamma_i,\alpha_i} |\psi_1|_0 |\psi_2|_0 |\rho|_0
$
for $i=1,2$, we obtain
$$
A(\gamma(t),\alpha(t)) \le C_{\gamma_1,\alpha_1}^{t}
C_{\gamma_2,\alpha_2}^{1-t} |\psi_1|_0 |\psi_2|_0 |\rho|_0
$$
thus proving eq.~(\ref{eq:dual-bound-1}) for the couple
$(\gamma(t),\alpha(t))$ for any $t\in [0,1]$. This implies that the
bound~(\ref{eq:dual-bound-1}) holds for any couple $(\gamma,\alpha)$
in the convex hull of the set $\DD_1 \cup \DD_2$. 
It is not difficult to see that this convex hull corresponds to the
region $\DD$ of the statement. This conclude the proof. 
\end{proof}



\subsection{Proof of Lemma~\ref{lemma:X2}}

\begin{proof}[]
We start the argument  as in the proof of Lemma~\ref{lemma:X}.
Let $\Psi_{ts}=X_{ts}(\varphi_1,\varphi_{21},\varphi_{22})$, then
$$
\Psi_{ts}(k) = - 2 \sum'_{k_1,k_{21}}  \frac{k k_2}4  \int_s^t  \int_s^\sigma e^{-i 3
  k k_1 k_2 \sigma
- i 3 k_2 k_{21} k_{22} \sigma_1}
 \, d\sigma   d\sigma_1\,  \varphi_1(k_1)   \varphi_{21}(k_{21})  \varphi_{22}(k_{22}).
$$  
Note that it is enough to restrict the sums over the set of $k_i$ for  which $k k_1 k_2 \neq 0$ and $k_2 k_{21} k_{22}\neq 0$. Then
the multiplier
\begin{equation}
  \label{eq:M-mult}
M = -   \frac{k k_2}2  \int_s^t  \int_s^\sigma e^{-i 3
  k k_1 k_2 \sigma
- i 3 k_2 k_{21} k_{22} \sigma_1} \, d\sigma   d\sigma_1;  
\end{equation}
has two different behaviors depending on the quantity $h = k k_1 k_2 + k_2 k_{21} k_{22}$ being zero or not. 
Split $\Psi$ accordingly and let
$$
\Psi^1_{ts}(k) = - 2 \sum'_{k_1,k_{21}: h= 0} M \,  \varphi_1(k_1)   \varphi_{21}(k_{21})  \varphi_{22}(k_{22})  
$$  
and $\Psi^2 = \Psi - \Psi^1$. We have $h=k_2(k_1+k_{21})(k-k_{21})$ so in the sum for $\Psi^1$ the only relevant contributions came from the case where $k_1+k_{21}=0$ or $k-k_{21}=0$. In the first case we have $k_{22}=k$ while in the second $k_{22}=-k_1$ so defining
$$
\Phi_{ts}(\varphi_1,\varphi_{21})(k) = - 2 \sum'_{k_1} M|_{k_{21}=-k_1,k_{22}=k} \,  \varphi_1(k_1)   \varphi_{21}(-k_{1})  
$$
we have $\Psi_{ts} = \Psi^{2}_{ts} + \Phi_{ts}(\varphi_1,\varphi_{21}) \varphi_{22} + \Phi_{ts}(\varphi_1,\varphi_{22}) \varphi_{21} $.

In Lemma~\ref{lemma:M-bounds} below we prove that, if $\gamma \le 1/2$ and $h \neq 0$ then
$
|M| \le C (A_1+A_2)  |t-s|^{2\gamma} 
$
where 
$$
A_1 = |k| |k_2| |k k_1 k_2|^{\gamma-1} |k_2 k_{21} k_{22}|^{\gamma-1}
$$
and
$$
A_2 = |k|^{1-\gamma} |k_2|^{2\gamma-1} |k_1|^{-\gamma}
|k-k_{21}|^{2\gamma-1} 
|k_1+k_{21}|^{2\gamma-1} |k_{21}|^{\gamma-1} |k_{22}|^{\gamma-1}
$$
while if $h=0$,
$
|M| \le C |k_1|^{-1}  |t-s| 
$.

\textbf{Bound for $\Phi$.}
We have
 $|\Phi(\varphi_1,\varphi_{21})| \le C |t-s| \Xi_0$ where
\begin{equation*}
\Xi_0 =  \sum_{k_1} |k_1|^{-2\alpha-1} |\rho(k)| |\psi_1(k_1)| |\psi_{21}(-k_1)|
\end{equation*}
It is clear that these expressions can be controlled with the $L^2$ norms of the test functions as soon as $-2\alpha-1 \le 0$.

\textbf{Bound for $\Psi^2$.}
In order to bound $\Psi^2$ it will be enough to bound separately
the forms
$$
\Xi_1 = \sum_{k,k_1,k_{21}} A_1 |k|^{\alpha} |k_1|^{-\alpha}
|k_{21}|^{-\alpha} |k_{22}|^{-\alpha} |\rho(k)| |\psi_1 (k_1)|   |\psi_{21}(k_{21})|  |\psi_{22}(k_{22})|  
$$
and
$$
\Xi_2 = \sum_{k,k_1,k_{21}} A_2 |k|^{\alpha} |k_1|^{-\alpha}
|k_{21}|^{-\alpha} |k_{22}|^{-\alpha} |\rho(k)| |\psi_1(k_1)|   |\psi_{21}(k_{21})|  |\psi_{22}(k_{22})|  
$$
in terms of the $L^2$ norms of the functions $\rho,\psi_1,\psi_{21},\psi_{22}$.

The expression $A_1$ is nicely factorized in its dependence on the
couple $k_{21},k_{22}$ and  the form $\Xi_1$ can be easily bounded
using (twice and iteratively) the same arguments as in
Lemma~\ref{lemma:X}.
We will concentrate on the form $\Xi_2$ which require  different
estimates. 

Write $\Xi_2$ as
$$
\Xi_2 = \sum_{k,k_1,k_{21}} |k|^{1-\gamma+\alpha} A_3
 |\rho(k)| |\psi_1(k_1)|   |\psi_{21}(k_{21})|  |\psi_{22}(k_{22})|  
$$
where
$
A_3 =    |k_1|^{-\gamma-\alpha}
|k_2(k-k_{21})(k_1+k_{21})|^{2\gamma-1} 
 |k_{21} k_{22}|^{\gamma-\alpha-1}
$
and use Cauchy-Schwartz on the double sum over $k_1,k_{21}$ to bound
$$
\Xi_2 \le \sum_{k} |k|^{1-\gamma+\alpha} 
 |\rho(k)|   
 \left[\sum_{k_1,k_{21}}A_3^2\right]^{1/2}
\left[\sum_{k_1,k_{21}} |\psi_1(k_1)|^2   |\psi_{21}(k_{21})|^2  |\psi_{22}(k_{22})|^2\right]^{1/2}
$$
Now Lemma~\ref{lemma:sum-bound} below states that for $(\gamma,\alpha) \in \DD$ we
have
$
\sup_{k\in\ZZ_*} |k|^{1-\gamma+\alpha} 
 [\sum_{k_1,k_{21}}A_3^2]^{1/2} \le C
$
so that
$$
\Xi_2 \le C \sum_{k}
 |\rho(k)|   
\left[\sum_{k_1,k_{21}} |\psi_1(k_1)|^2   |\psi_{21}(k_{21})|^2  |\psi_{22}(k_{22})|^2\right]^{1/2}.
$$
Then using again  Cauchy-Schwartz (on the $k$ summation) we end up with
$$
\Xi_2 \le C |\rho|_0
\left[\sum_{k,k_1,k_{21}} |\psi_1(k_1)|^2   |\psi_{21}(k_{21})|^2
  |\psi_{22}(k_{22})|^2\right]^{1/2} = C|\rho|_0 |\psi_1|_0
|\psi_{21}|_0 |\psi_{22}|_0
$$
where the last equality is obtained after a simple change of variables
in the summations. This ends the proof of the statement.
\end{proof}

\bigskip
It remains to prove some complementary results. We will use the
notation introduced above in the proof of Lemma~\ref{lemma:X2}.

\begin{lemma}
\label{lemma:M-bounds}
If $\gamma \in[0,1/2]$ and $h \neq 0$ we have
$$
|M| \le C (A_1+A_2)  |t-s|^{2\gamma} 
$$  
If $h=0$
$$
|M| \le C |k_1|^{-1}  |t-s|
$$  
\end{lemma}
\begin{proof}
  Take two numbers $a,b \in \RR$ and let 
$
m = \int_s^t \int_s^\sigma e^{i a \sigma} e^{i b \sigma_1} d\sigma d\sigma_1 
$.
An explicit integration gives
$$
m = \frac{e^{i(a+b)s}-e^{i(a+b)t}}{b (a+b)} 
+\frac{e^{i bs+iat}-e^{i(a+b)s}}{ab}.
$$
Then
\begin{equation}
  \label{eq:m-bound-1}
|m| \le \frac{1}{|b||a+b|}+\frac{1}{|a||b|}
\end{equation}
moreover
\begin{equation*}
|m| \le  \int_s^t \left|\int_s^\sigma  e^{i b \sigma_1}
  d\sigma_1\right| d\sigma \le \frac{|t-s|}{|b|}  
\end{equation*}
and symmetrically $|m| \le |t-s| |a|^{-1}$. These last bounds imply
that  $|m| \le |t-s||ab|^{-1/2}$ and moreover imply the bound in the case $h=0$.
Otherwise, since $2\gamma \in [0,1]$,  interpolating between this bound and
eq.~(\ref{eq:m-bound-1}) we get
\begin{equation}
  \label{eq:m-bound-3}
 |m| \le \frac{|t-s|^{2\gamma}}{|ab|^{\gamma}}
 \left(\frac{1}{|b||a+b|}+\frac{1}{|a||b|}\right)^{1-2\gamma} \le 
\frac{|t-s|^{2\gamma}}{|b|^{\gamma}|a|^{1-\gamma}|a+b|^{1-2\gamma}}+\frac{|t-s|^{2\gamma}}{|a|^{1-\gamma}|b|^{1-\gamma}}
\end{equation}
which, setting $a = 3k k_1 k_2$ and $b = 3 k_2 k_{21} k_{22}$ implies
the statement we had to prove.
\end{proof}


\begin{lemma}
\label{lemma:sum-bound}
Let
$$
S_\eps(k) =  \sum_{k_1,k_{21}}   |k_2|^{4\gamma-2+2\eps} |k_1|^{-2\gamma-2\alpha+2\eps}
|(k-k_{21})(k_1+k_{21})|^{4\gamma-2} 
 |k_{21} k_{22}|^{2\gamma-2\alpha+2\eps-2}
$$
Then if $(\gamma,\alpha) \in \DD$, there exists $\eps > 0$ such that
$$
 |k|^{2-2\gamma+2\alpha+\eps} S_{\eps}(k)
 \le C
$$  
uniformly on $k \in \ZZ_*$.
\end{lemma}
\begin{proof}
We will prove the statement for $\eps = 0$, the extension for a positive but small $\eps$ being a straightforward modification of the resulting argument.

First we bound the double sum by a planar integral. Define the integration region $\cR(z,a)$ with $a>0$, $z\in\RR$ as the set of points $(x,y) \in \RR^2$ such that
$$
|x|>a, |y| > a, |x-z|>a, |y-z|>a, |x+y|>a, |x+y-z|>a
$$
Then 
$$
S_0(k) \le C \int_{\cR(k,1)} dx dy |k-x|^{4\gamma-2} |x|^{-2\gamma-2\alpha}
|(k-y)(x+y)|^{4\gamma-2} 
 |y (k-x-y)|^{2\gamma-2\alpha-2}
$$
which, by scaling away $|k|$ and by reflection symmetry, is equivalent to
$$
S_0(k) \le C |k|^{2(4\gamma-2)+3(2\gamma-2\alpha-2)+2} I(\cR(1,|k|^{-1}))
$$
with
$$
I(\mathcal{A}) =
\int_{\mathcal{A}} dx dy |1-x|^{4\gamma-2} |x|^{-2\gamma-2\alpha}
|(1-y)(x+y)|^{4\gamma-2} 
 |y (1-x-y)|^{2\gamma-2\alpha-2}.
$$
To prove the original statement we must show that
$
|a|^{2(3-6\gamma+2\alpha)} I(\cR(1,a)) \le C
$ uniformly in $a \to 0$.

We start by showing that $I(\cR(1,a))$ is actually finite for any $a > 0$. Then we will find estimates for $I(a)$ as $a \to 0$.

\paragraph{Boundedness}
Note that, when $|x|+|y|$ is large we can bound the integrand of $I(\cR(1,a))$ by
$$
C |x|^{2\gamma-2\alpha-2}
|y(x+y)|^{6\gamma-2\alpha-24}
$$
Then since in $\cR(1,a)$ we have $|x|>a$ we can bound this last integrand by
$C |a|^{2\gamma-2\alpha-2}
|y(x+y)|^{6\gamma-2\alpha-4}$. Using the fact  that $6\gamma-2\alpha-4 < -1$ in $\DD$, the integral
$
\int_{\cR(1,a)} |y(x+y)|^{6\gamma-2\alpha-4} dx dy
$ can be shown to be finite.

\paragraph{Asymptotics} When $a \to 0$ the integral can blow up due to contributions from the sub-domain $\cR'(a) = \cR(1,a) \backslash \cR(1,a')$ where $a'>0$ is a small but fixed quantity.  Denote with $B(z,r) \subset \RR^2$ the ball  centered in $z$ and of radius $r$. The region $\cR'(a)$ can be covered by the union of the  seven sets 
$$
\cB_i = \cR'(a) \cap B(z_i, a'),\,\, i=A,\dots,F\,\,, \qquad \cR''(a) = \cR'(a)\backslash \left[\cup_{i=A,\dots,F} \cB_i \right]
$$
where $z_i$,$i=A,\dots,F$ are the following points
$$
z_A = (-1,1),\, z_B = (0,1),\, z_C=(1,1),\, z_D=(0,0),\, z_E=(1,0),\, z_F=(1,-1).
$$
We will study the contributions to $I(\cR'(a))$ due to $I(\cB_i)$, $i=A,\dots,F$. It is not difficult then to show that $I(\cR''(a))$ gives smaller contributions to the $a\to 0$ limit.

Let us show how to handle $I(\cB_A)$. First bound
$$
I(\cB_A) \le C \int_{\cB_A} 
|(1-y)(x+y)|^{4\gamma-2} d x\, d y
$$
the other factors being bounded as $a \to 0$ for $(x,y) \in \cB_A$. By a change of variables we can reduce this integral to the following  form (eventually increasing the integration domain)
$
I(\cB_A) \le C \int_{a}^{a''} |x y|^{4\gamma-2} dx dy
$
where $a''>a$ is another fixed constant. Then we easily compute the factorized integral and finally  bound the original expression with
$$
I(\cB_A) \le C a^{-2(1-4\gamma)_+} = C a^{\zeta_A}
$$
where $(\cdot)_+$ denotes the positive part. The constant $C$ in this expression does not depend on $a$. 

The treatment of all $I(\cB_i)$ for $i=B,\dots,F$ follows similar lines. To simplify the argument it is not difficult to see that $I(\cB_i)$ for $i=A,C,F$  have the same asymptotic bound. The same is true for $I(\cB_i)$ for $i=B,D$ where
$$
I(\cB_{B,D}) \le C \int_{\cB_B} 
|x|^{-2\gamma-2\alpha}|(1-y)|^{4\gamma-2} |x+y-1|^{2\gamma-2\alpha-2} d x\, d y
$$
Moreover
$$
I(\cB_{E}) \le C\int_{\cB_E} |x-1|^{4\gamma-2} |y|^{2\gamma-2\alpha-2} |x+y-1|^{2\gamma-2\alpha-2} dx\, dy
$$
and since $-2\gamma-2\alpha > 2\gamma-2\alpha-2$ for $\gamma \in [0,1/2)$ we have also
$
I(\cB_{B,D}) \le C I(\cB_E)
$.
Since we already have established the behavior of $I(\cB_A)$ it remains only to understand that of $I(\cB_E)$. We have two possibilities.
The first is to bound $|x+y-1|^{2\gamma-2\alpha-2} \le |a|^{2\gamma-2\alpha-2}$ in $\cB_E$, then
\begin{equation*}
  \begin{split}
I(\cB_{E}) & \le C |a|^{2\gamma-2\alpha-2} \int_{\cB_E} |x-1|^{4\gamma-2} |y|^{2\gamma-2\alpha-2}  dx\, dy
\\ & \le C |a|^{2\gamma-2\alpha-2 - (1-4\gamma)_+ - (1-2\gamma+2\alpha)_+}  = C a^{\zeta_{E,1}}  
  \end{split}
\end{equation*}
by reasoning as in the case of $I(\cB_A)$ above. The second is to bound $|x-1|^{4\gamma-2}\le |a|^{4\gamma-2}$ to obtain
\begin{equation*}
  \begin{split}
I(\cB_{E}) & \le C |a|^{4\gamma-2} \int_{\cB_E} |x+y-1|^{2\gamma-2\alpha-2} |y|^{2\gamma-2\alpha-2}  dx\, dy
\\ & \le C |a|^{4\gamma-2 - 2(1-2\gamma+2\alpha)_+}  = C a^{\zeta_{E,2}}  
  \end{split}
\end{equation*}
and let $\zeta_E = \max(\zeta_{E,1},\zeta_{E,2})$.

In the same way it is not difficult to show that $I(\cR''(a)) \le C (|a|^{\zeta_A}+ |a|^{\zeta_E})$. Putting all these estimates together we get
$
I(\cR(1,a)) \le C (|a|^{\zeta_A}+ |a|^{\zeta_E})
$. It remains to prove that
\begin{equation}
  \label{eq:cases-res}
\zeta'_{A,E} = 2(3-6\gamma+2\alpha)+\zeta_{A,E}  \ge 0  
\end{equation}
for any admissible couple $(\gamma,\alpha)$.
We proceed by considering three different regions $C_{1,2,3}$ of admissible couples:
\begin{itemize}
\item[$C_1$]: $\gamma \in [0,1/4)$, $1-2\gamma+2\alpha < 0$. Here $\zeta_A = -2(1-4\gamma)$ and $\zeta_E = \zeta_{E,1} = 6\gamma-2\alpha-3$. Then $\zeta'_A = 2(2-2\gamma+2\alpha)$ and $\zeta'_E = 3-6\gamma+2\alpha$;

\item[$C_2$]: $\gamma \in [0,1/4)$, $1-2\gamma+2\alpha > 0$. Here $\zeta_A = -2(1-4\gamma)$ and $\zeta_E= \zeta_{E,1} = 8\gamma-4\alpha-4$. Here $\zeta'_A = 2(2-2\gamma+2\alpha)$ and $\zeta'_E = 2(1-2\gamma)$;
 
\item[$C_3$]: $\gamma \in [1/4,1/2)$, $1-2\gamma+2\alpha > 0$. Here $\zeta_A = 0$ and $\zeta_E = \zeta_{E,2} =  4\gamma-4\alpha-3$. Here $\zeta'_A = 2(3-6\gamma+2\alpha)$ and $\zeta'_E = 2(1-2\gamma)$.

\end{itemize}
We can check that in the three regions $C_{1,2,3}$ and for $(\gamma,\alpha) \in \DD$ we always end up with $\zeta'_{A,E} > 0$. By interpolation between these three regions we  prove the claim for couples in their convex hull which coincide with $\DD$.

To treat the case of small positive $\eps$ we just retrace this argument and take $\eps$ small enough so that all relevant inequalities are still satisfied. 
\end{proof}



\subsection{Proof of Lemma~\ref{lemma:gal}}


\begin{proof}
Let us work out the details for $X$, first.
Since $(\gamma,\alpha) \in \DD$, by Lemma~\ref{lemma:X} there exists $\eps > 0$ such that $X \in \CC_2^\gamma \LL((H^{\alpha-\eps})^2,H^{\alpha+\eps})$.
By definition
$$
|X^{(N)}_{ts}(\varphi_1,\varphi_2)|_{\alpha} = | P_N  X_{ts}(P_N \varphi_1, P_n \varphi_2)|_\alpha
$$  
Then
\begin{equation}
\label{eq:gal1}
  \begin{split}
 | (X^{(N)}_{ts}-X_{ts})(\varphi_1,\varphi_2)|_{\alpha} & \le 
|\wt P_N X_{ts}(P_N \varphi_1,P_N \varphi_2)|_\alpha + |X_{ts}(\wt P_N \varphi_1,P_N \varphi_2)|_\alpha 
\\ & \qquad + |X_{ts}( \varphi_1,\wt P_N \varphi_2)|_\alpha 
  \end{split}
\end{equation}
where $\wt P_N = \text{Id}-P_N$ is the projector over modes with $|k| > N$. Now note that, for $\varphi \in H^{\alpha}$ we have 
$
|\wt P_N \varphi|_{\alpha-\eps} \le N^{-\eps} |\varphi|_\alpha
$
so the three terms in eq.~(\ref{eq:gal1}) can be bounded as follows.
The first as
\begin{equation*}
  \begin{split}
|\wt P_N X_{ts}(P_N \varphi_1,P_N \varphi_2)|_\alpha & \le N^{-\eps} |X_{ts}(P_N \varphi_1,P_N \varphi_2)|_{\alpha+\eps} 
\\ & \le C N^{-\eps} |P_N \varphi_1|_{\alpha-\eps} |P_N \varphi_2|_{\alpha-\eps}   
\le C N^{-\eps} |\varphi_1|_{\alpha} |\varphi_2|_{\alpha}     ,  
  \end{split}
\end{equation*}
the second as
\begin{equation*}
  \begin{split}
|X_{ts}(\wt P_N \varphi_1,P_N \varphi_2)|_\alpha  \le C  |\wt P_N \varphi_1|_{\alpha-\eps} |P_N \varphi_2|_{\alpha-\eps} 
 \le C  N^{-\eps} |\varphi_1|_{\alpha} |\varphi_2|_{\alpha} ,
  \end{split}
\end{equation*}
 and a similar bound for the third.
Then
\begin{equation}
\label{eq:gal2}
  \begin{split}
 | (X^{(N)}_{ts}-X_{ts})(\varphi_1,\varphi_2)|_{\alpha} & \le C N^{-\eps} |\varphi_1|_{\alpha} |\varphi_2|_{\alpha}
  \end{split}
\end{equation}
which implies the claimed convergence.

Now consider the convergence of $\hat X^{(N),2}$ to $\hat X^{(2)}$. We want to show that
\begin{equation}
\label{eq:gal-X2-1}
  \begin{split}
 | (\hat X^{(N),2}_{ts}-\hat X^2_{ts})(\varphi_1,\varphi_2,\varphi_3)|_{\alpha} & \le C N^{-\eps} |\varphi_1|_{\alpha} |\varphi_2|_{\alpha} |\varphi_3|_{\alpha}
  \end{split}
\end{equation}
for some $\eps > 0$, this will imply the claim. 
Define a multiplier $\Theta$ as
$$
\Theta =  \chi_N(k) \chi_N(k_1) \chi_N(k_2) \chi_N(k_{21}) \chi_N(k_{22})-1  
$$
where $\chi_N(k) = 1_{|k|\le N}$ is the kernel in the Fourier basis of the  operator $P_N$.
By duality we need to bound the form
\begin{equation}
\label{eq:form-bound}
  \begin{split}
\sum_{k,k_1,k_{21}} |M| |\Theta| |k|^{\alpha} |k_1|^{-\alpha}
|k_{21}|^{-\alpha} |k_{22}|^{-\alpha} |\rho(k)| |\psi_1(k_1)|   |\psi_{21}(k_{21})|  |\psi_{22}(k_{22})|  
  \end{split}
\end{equation}
with $C N^{-\eps} |t-s|^{2\gamma} |\rho|_0 |\psi_1|_0 |\psi_{21}|_0 |\psi_{22}|_0$. Recall that the multiplier $M$ has been defined in eq.~(\ref{eq:M-mult}). 
The $\Theta$ multiplier can be rewritten as a telescopic sum
\begin{equation*}
  \begin{split}
\Theta & = [\chi_N(k)-1] \chi_N(k_1) \chi_N(k_2) \chi_N(k_{21}) \chi_N(k_{22}) +
[\chi_N(k_1)-1] \chi_N(k_2) \chi_N(k_{21}) \chi_N(k_{22})
\\ & +  [\chi_N(k_2)-1] \chi_N(k_{21}) \chi_N(k_{22}) +  [\chi_N(k_{21})-1] \chi_N(k_{22}) + [\chi_N(k_{22})-1]    
  \end{split}
\end{equation*}
and then can be bounded as
$$
|\Theta| \le |\chi_N(k)-1|+|\chi_N(k_1)-1|+|\chi_N(k_2)-1|+|\chi_N(k_{21})-1|+|\chi_N(k_{22})-1|
$$
Moreover since
$$
|\chi_N(k)-1| \le |\chi_N(k)-1| \frac{|k|^\eps}{N^\eps} \le N^{-\eps} |k|^\eps 
$$
for any $\eps \ge 0$ we have that
$$
|\Theta| \le N^{-\eps}(|k|^\eps+|k_{1}|^\eps+|k_{2}|^\eps+|k_{21}|^\eps+|k_{22}|^\eps)
\le N^{-\eps} |k k_1 k_2 k_{21} k_{22}|^\eps
$$
Then to obtain the required $N$ dependence of the bound for the form~(\ref{eq:form-bound}) it is enough to bound the new $N$-independent  form
\begin{equation}
\label{eq:form-bound-new}
  \begin{split}
\sum_{k,k_1,k_{21}} |M|  |k|^{\alpha+\eps} |k_1|^{-\alpha+\eps} |k_2|^{\eps}
|k_{21}|^{-\alpha+\eps} |k_{22}|^{-\alpha+\eps} |\rho(k)| |\psi_1(k_1)|   |\psi_{21}(k_{21})|  |\psi_{22}(k_{22})|  
  \end{split}
\end{equation}
with a quantity like $C  |t-s|^{2\gamma} |\rho|_0 |\psi_1|_0 |\psi_{21}|_0 |\psi_{22}|_0$. By retracing the proof of Lemma~\ref{lemma:X2} and using the $\eps>0$ version of Lemma~\ref{lemma:sum-bound} we can prove that such bound holds for all admissible couples $(\gamma,\alpha)$. 

This completes the proof of the statement.

\end{proof}


%\bibliographystyle{plain}
%\bibliography{kdv}


\begin{thebibliography}{10}

\bibitem{evolution}
Hakima Bessaih, Massimiliano Gubinelli, and Francesco Russo.
\newblock The evolution of a random vortex filament.
\newblock {\em Ann. Probab.}, 33(5):1825--1855, 2005.

\bibitem{MR1215780}
J.~Bourgain.
\newblock Fourier transform restriction phenomena for certain lattice subsets
  and applications to nonlinear evolution equations. {II}. {T}he
  {K}d{V}-equation.
\newblock {\em Geom. Funct. Anal.}, 3(3):209--262, 1993.

\bibitem{MR2018661}
Michael Christ, James Colliander, and Terrence Tao.
\newblock Asymptotics, frequency modulation, and low regularity ill-posedness
  for canonical defocusing equations.
\newblock {\em Amer. J. Math.}, 125(6):1235--1293, 2003.

\bibitem{MR1969209}
J.~Colliander, M.~Keel, G.~Staffilani, H.~Takaoka, and T.~Tao.
\newblock Sharp global well-posedness for {K}d{V} and modified {K}d{V} on
  {$\mathbb{R}$} and {$\mathbb{T}$}.
\newblock {\em J. Amer. Math. Soc.}, 16(3):705--749 (electronic), 2003.

\bibitem{MR2164030}
Laure Coutin and Antoine Lejay.
\newblock Semi-martingales and rough paths theory.
\newblock {\em Electron. J. Probab.}, 10:no. 23, 761--785 (electronic), 2005.

\bibitem{MR1883719}
Laure Coutin and Zhongmin Qian.
\newblock Stochastic analysis, rough path analysis and fractional {B}rownian
  motions.
\newblock {\em Probab. Theory Related Fields}, 122(1):108--140, 2002.

\bibitem{davie}
A.~M. Davie.
\newblock Differential equations driven by rough signals: an approach via
  discrete approximation.
\newblock preprint, 2003.

\bibitem{MR2111917}
A.~De~Bouard, A.~Debussche, and Y.~Tsutsumi.
\newblock Periodic solutions of the {K}orteweg-de {V}ries equation driven by
  white noise.
\newblock {\em SIAM J. Math. Anal.}, 36(3):815--855 (electronic), 2004/05.

\bibitem{fvgeo}
P.~Friz and N.~Victoir.
\newblock A note on the notion of geometric rough paths.
\newblock to appear in Prob. Th. Rel. Fields, 2004.

\bibitem{fveuler}
P.~Friz and N.~Victoir.
\newblock Euler estimates for rough differential equations.
\newblock preprint, 2006.

\bibitem{fvsubell}
P.~Friz and N.~Victoir.
\newblock On uniformly subelliptic operators and stochastic area.
\newblock preprint, 2006.

\bibitem{MR1470933}
J.~G. Gaines and T.~J. Lyons.
\newblock Variable step size control in the numerical solution of stochastic
  differential equations.
\newblock {\em SIAM J. Appl. Math.}, 57(5):1455--1484, 1997.

\bibitem{MR1019307}
J.~Ginibre and Y.~Tsutsumi.
\newblock Uniqueness of solutions for the generalized {K}orteweg-de {V}ries
  equation.
\newblock {\em SIAM J. Math. Anal.}, 20(6):1388--1425, 1989.

\bibitem{MR1423623}
Jean Ginibre.
\newblock Le probl\`eme de {C}auchy pour des {EDP} semi-lin\'eaires
  p\'eriodiques en variables d'espace (d'apr\`es {B}ourgain).
\newblock {\em Ast\'erisque}, (237):Exp.\ No.\ 796, 4, 163--187, 1996.
\newblock S\'eminaire Bourbaki, Vol.\ 1994/95.

\bibitem{MR2091358}
M.~Gubinelli.
\newblock Controlling rough paths.
\newblock {\em J. Funct. Anal.}, 216(1):86--140, 2004.

\bibitem{trees}
M.~Gubinelli.
\newblock Branched rough paths.
\newblock in prepraration, 2006.

\bibitem{TindelGubinelli}
M.~Gubinelli and S.~Tindel.
\newblock Rough evolution equations.
\newblock 2006.

\bibitem{MR2181576}
T.~Kappeler and P.~Topalov.
\newblock Well-posedness of {K}d{V} on {$H\sp {-1}(\mathbb{T})$}.
\newblock In {\em Mathematisches Institut, Georg-August-Universit\"at
  G\"ottingen: Seminars 2003/2004}, pages 151--155. Universit\"atsdrucke
  G\"ottingen, G\"ottingen, 2004.

\bibitem{MR759907}
Tosio Kato.
\newblock On the {C}auchy problem for the (generalized) {K}orteweg-de {V}ries
  equation.
\newblock In {\em Studies in applied mathematics}, volume~8 of {\em Adv. Math.
  Suppl. Stud.}, pages 93--128. Academic Press, New York, 1983.

\bibitem{MR1086966}
Carlos~E. Kenig, Gustavo Ponce, and Luis Vega.
\newblock Well-posedness of the initial value problem for the {K}orteweg-de
  {V}ries equation.
\newblock {\em J. Amer. Math. Soc.}, 4(2):323--347, 1991.

\bibitem{MR1230283}
Carlos~E. Kenig, Gustavo Ponce, and Luis Vega.
\newblock The {C}auchy problem for the {K}orteweg-de {V}ries equation in
  {S}obolev spaces of negative indices.
\newblock {\em Duke Math. J.}, 71(1):1--21, 1993.

\bibitem{MR1329387}
Carlos~E. Kenig, Gustavo Ponce, and Luis Vega.
\newblock A bilinear estimate with applications to the {K}d{V} equation.
\newblock {\em J. Amer. Math. Soc.}, 9(2):573--603, 1996.

\bibitem{youngspde}
Antoine Lejay, Massimiliano Gubinelli, and Samy Tindel.
\newblock Young integrals and {SPDE}s.
\newblock {\em Pot. Anal.}, 2006.

\bibitem{MR2036784}
Terry Lyons and Zhongmin Qian.
\newblock {\em System control and rough paths}.
\newblock Oxford Mathematical Monographs. Oxford University Press, Oxford,
  2002.
\newblock Oxford Science Publications.

\bibitem{MR1654527}
Terry~J. Lyons.
\newblock Differential equations driven by rough signals.
\newblock {\em Rev. Mat. Iberoamericana}, 14(2):215--310, 1998.

\end{thebibliography}
\end{document}

%%% Local Variables: 
%%% mode: latex
%%% TeX-master: t
%%% End: 
